% =====================================================================================================================
% Experimental setup and results for "Successive Convexification In Learned World Models For Safe Offline Planning".
% Main text: % =====================================================================================================================
% Experimental setup and results for "Successive Convexification In Learned World Models For Safe Offline Planning".
% Main text: % =====================================================================================================================
% Experimental setup and results for "Successive Convexification In Learned World Models For Safe Offline Planning".
% Main text: % =====================================================================================================================
% Experimental setup and results for "Successive Convexification In Learned World Models For Safe Offline Planning".
% Main text: \input{experiments_results}. Appendix: \input{appendix_results}. Citations: experiments_refs.bib.
%
% Every figure and table is written by the repository's scripts, run from the repository root:
%   python src/scripts/analyse_results.py  -> src/save_data/experiments/analysis_results/*.csv
%   python src/scripts/paper_figures.py    -> paper_graphs/{main_text,ablation,appendix}/ (figures and both tables)
%   python src/scripts/env_figure.py       -> paper_graphs/main_text/fig_environments.pdf
% \figdir is the paper_graphs folder; define it before \input if the paper lives elsewhere (e.g. on Overleaf).
%
% Notation follows the method: N horizon, \epsilon barrier margin, \alpha DCBF rate, c Lipschitz bound, \tau risk
% threshold, g_\Phi margin head. Labels of the method's parts used here (rename to match yours): sec:regularity,
% sec:cbf, sec:scp, sec:init, eq:ocp.
% Every number is from the analysis of the runs on main (2d4240b). "% TODO" marks a fact to confirm before submission.
% =====================================================================================================================
\providecommand{\figdir}{../paper_graphs}

\section{Experimental Setup}\label{sec:setup}

Our experiments answer four questions.
\textbf{(Q1)}~Does SCP, started from a sampled plan, plan better than sampling-based MPC given the same wall-clock time?
\textbf{(Q2)}~How does the comparison change as the horizon, and with it the number of decision variables, grows?
\textbf{(Q3)}~What does SCP's convergence tell us online?
\textbf{(Q4)}~Does meeting the barrier in the learned model keep the robot safe in the environment, and what do the
regularity constraints cost?

% ---------------------------------------------------------------------------------------------------------------------
\subsection{Environments}\label{sec:envs}

\begin{figure}[t]
  \centering
  \includegraphics[width=0.62\linewidth]{\figdir/main_text/fig_environments.pdf}
  \caption{The two tasks from above, one layout each (layouts are drawn from the episode's seed). The point robot
  (red) must reach the goal (green) without entering a hazard (blue); in PointGoal2, touching a vase (cyan) is also a
  violation.}
  \label{fig:envs}
\end{figure}

We use \texttt{SafetyPointGoal1} and \texttt{SafetyPointGoal2} from Safety-Gymnasium~\citep{ji2023safety}
(Fig.~\ref{fig:envs}). A point robot with two actuators (forward force and turning, $u\in[-1,1]^2$) must reach a goal
of radius $0.3$\,m, which moves to a new position each time it is reached, within $T=1000$ steps. The reward is the
per-step reduction in distance to the goal plus a bonus of 1 for reaching it. The robot observes 60 quantities: 12
proprioceptive signals (accelerometer, velocimeter, gyroscope and magnetometer) and three 16-ray pseudo-lidars measuring
the proximity of the goal, the hazards and the vases. Level~1 places 8 hazards of radius $0.2$\,m and one cost-free
vase in a $3\times3$\,m region; level~2 places 10 hazards and 10 vases in a $4\times4$\,m region and also charges for
touching or moving a vase. A step is a \emph{violation} if it incurs any cost. We treat violations as absorbing
failures: an episode ends at its first violation, so the \emph{violation rate} is the share of episodes that fail and
the \emph{return} is the reward collected before failure or the time limit.

These tasks suit our question for four reasons.
(i)~DSRL~\citep{liu2023datasets} provides offline datasets for both, so the world model is learned without any
interaction, which is the setting of this paper.
(ii)~The constraint is geometric (stay out of regions of the plane) and observed through the lidar, so the barrier's
margin is a distance in metres, the margin $\epsilon$ has a physical meaning, and safety in the model can be checked
against safety in the environment (Sec.~\ref{sec:latent_safety}).
(iii)~Level~2 is more cluttered (20 obstacles in 16\,m$^2$ against 8 in 9\,m$^2$) and adds a second, movable kind of
obstacle, so the two levels, together with $\epsilon$, vary how restrictive the safe set is.
(iv)~The action is two-dimensional, which favours sampling: at $N=10$ a plan has only 20 decision variables, so the
comparison is a demanding one for SCP. The horizon study (Sec.~\ref{sec:horizon}) raises this to 40 and 60.

% ---------------------------------------------------------------------------------------------------------------------
\subsection{Offline data and world models}\label{sec:wm_setup}

For each task we train on the corresponding DSRL dataset (2,022 trajectories for PointGoal1 and 3,442 for PointGoal2,
of up to 1000 steps each), truncating every trajectory at its first violation to match the absorbing failure above.
This leaves about 0.87\,M transitions for PointGoal1 and 1.15\,M for PointGoal2; only 329 of PointGoal2's trajectories
run their full length without a violation.
% TODO: PG1's exact transition count (FilterDataset.ipynb prints it for PG2 only: 1,145,452); one sentence on how DSRL
% collected the data (behaviour policies, cost budgets).
The encoder drops seven proprioceptive channels (the accelerometer and the out-of-plane components of velocity,
angular rate and magnetic field) and maps the remaining 53 to a latent state of $d_z=32$ dimensions. The four networks
the planner differentiates through (dynamics, reward, value and margin) are MLPs with two hidden layers of 256 units,
each with Lipschitz constant at most $c$ (Sec.~\ref{sec:regularity}). We train one model per task and bound,
$c\in\{1.2,1.5,\infty\}$ ($c=\infty$: unbounded), for $N=10$, and two more PointGoal1 models with $c=1.5$ for $N=20$
and $30$. Each is trained for 250 epochs of 1000 (PointGoal1) or 1500 (PointGoal2) minibatches of 256 windows of $N+1$
consecutive transitions, with learning rate $3\times10^{-4}$, the loss at rollout step $t$ weighted by $\rho^t$
($\rho=0.9$, $0.95$, $0.99$ for $N=10$, $20$, $30$), discount $\gamma=0.99$ and value expectile $0.7$. The margin head
$g_\Phi$ regresses the distance from the robot to the nearest hazard edge (in PointGoal2, hazard or vase edge),
computed from the lidar, clipped at 1\,m, with samples near an edge up-weighted.
% TODO: optimiser (Adam?); state that each model is a single training seed (limitations).

% ---------------------------------------------------------------------------------------------------------------------
\subsection{Planners}\label{sec:planners}

All planners solve the same receding-horizon problem~\eqref{eq:ocp} in the same model: maximise
$\sum_{k=0}^{N-1}\gamma^k r(z_k,u_k)+\gamma^N V(z_N)-\rho_u\|u\|^2$ ($\rho_u=10^{-4}$) subject to
$(1-\alpha)h(z_k)-h(z_{k+1})\le\tau=0$ at every step, with $h(z)=g_\Phi(z)-\epsilon$ and $\alpha=0.05$. Each executes
the first action of its plan and replans at every step.

\textbf{SCP} is the method of Sec.~\ref{sec:scp} started from its previous plan shifted by one step (zeros at the first
step): at most 40 prox-linear iterations, each a convex QP solved by Clarabel~\citep{goulart2024clarabel}, with exact
$\ell_1$ penalties $w_{\mathrm{dyn}}=10^3$ and $w_{\mathrm{safe}}=10^2$, a trust-region weight that starts at 1, is
floored at $10^{-2}$ and is carried between steps, acceptance thresholds $(0.25,0.75)$, and a stop when the
prox-gradient norm falls below $10^{-2}$.
\textbf{SCP+MPPI}, our method, first runs MPPI with 512 sequences (below) and starts SCP from whichever of MPPI's plan
and its own shifted plan has the lower penalty merit (Sec.~\ref{sec:init}). SCP accepts only steps that lower the
merit, so the plan SCP+MPPI executes is never worse under the merit than MPPI's, whether or not SCP converges.
\textbf{MPPI}~\citep{williams2017information} and \textbf{CEM}~\citep{rubinstein1999cross} run 6 iterations of $J$
action sequences drawn from a Gaussian that is refitted to the best $J/8$ (weighted by $\exp(\text{score}/0.5)$ in MPPI,
uniformly in CEM) and warm-started from the previous mean shifted by one step. A sequence whose rollout breaks the
barrier at any step ranks below every sequence that meets it; if none meets it, sequences are ranked by their total
barrier excess, the quantity SCP's $\ell_1$ penalty charges. The executed plan is the refitted mean if its rollout
meets the barrier, and otherwise the best sequence sampled.

We run the samplers with $J=512$, the proposal SCP+MPPI uses, and \emph{time-matched}, with $J$ chosen so that a step
takes about as long as SCP+MPPI's: $J=65{,}536$ at $N=10$ (about 0.70\,s against 0.40--0.71\,s for SCP+MPPI, timed on
the same idle GPU), and $95{,}040$ and $120{,}000$ at $N=20$ and $30$ (mean logged plan times 2.04\,s and 3.81\,s,
against 2.16\,s and 3.61\,s for SCP+MPPI). Each run used one 1g.10gb slice of an NVIDIA A100, with the QP solved on the
CPU.
% TODO: the 65,536 time is interpolated between 32,768 (0.33 s) and 131,072 (1.43 s) sequences; time it directly or
% say so. Confirm the hardware sentence.

% ---------------------------------------------------------------------------------------------------------------------
\subsection{Protocol and statistics}\label{sec:protocol}

We sweep the margin $\epsilon\in\{0,0.1,0.2,0.3,0.5\}$\,m on every $N=10$ model and fix $\epsilon=0.2$\,m for the
horizon study. Each configuration runs on seeds $0$--$39$, or $0$--$29$ for the time-matched samplers, the horizon
study, PointGoal2's $c=\infty$ model and SCP+MPPI on PointGoal1's $c=\infty$ model. The seed fixes the layout, the
initial state and the planner's random numbers, so every comparison between planners or models is paired by seed. We
report violation rates with Wilson 95\% intervals and mean returns with bootstrap 95\% intervals ($10^4$ resamples).
Paired comparisons report the difference in violation rate and in return with a 95\% bootstrap interval, and the
probability, given the data, that SCP+MPPI's violation rate is lower or its return higher than the baseline's, from a
Bayesian bootstrap~\citep{rubin1981bayesian}. When pooling over margins and tasks we use $\epsilon\ge0.1$, since at
$\epsilon=0$ every planner fails often (Sec.~\ref{sec:q1}), and resample whole seeds, since a seed's episodes share a
layout. The per-model results in the appendix use McNemar's exact test on violations, Holm-corrected over the six models,
and the Wilcoxon signed-rank test on returns.

% =====================================================================================================================
\section{Results}\label{sec:results}

% ---------------------------------------------------------------------------------------------------------------------
\subsection{SCP needs a sampled start, and then matches time-matched sampling at \texorpdfstring{$N=10$}{N=10}}\label{sec:q1}

\begin{figure}[t]
  \centering
  \includegraphics[width=\linewidth]{\figdir/main_text/fig_eps_sweep_c1_5.pdf}
  \caption{Violation rate (top) and return (bottom) against the margin $\epsilon$ on each task's $c=1.5$ model at
  $N=10$, with 95\% intervals; 40 episodes per point (30 for the time-matched samplers).}
  \label{fig:eps_sweep}
\end{figure}

\begin{table}[t]
  \centering
  \caption{SCP+MPPI minus each baseline on the $c=1.5$ world models at $N=10$, pooled over both tasks and
  $\epsilon\ge0.1$ and paired by task, margin and seed (Tables~\ref{tab:pooled_c1_2} and~\ref{tab:pooled_cinf} in the
  appendix give the other bounds). A negative $\Delta$ violation rate or a positive $\Delta$ return favours SCP+MPPI.
  Brackets: 95\% intervals resampling whole seeds. $P$(safer) and $P$(higher return): the probability, given the data,
  that SCP+MPPI's violation rate is lower, or its return higher, than the baseline's on these models (Bayesian bootstrap
  over seeds).}
  \label{tab:pooled}
  \small
  \input{\figdir/main_text/table_pooled_c1_5}
\end{table}

\textbf{SCP needs a sampled start.} Started from its own shifted plan, SCP fails in 55\% of episodes over all $N=10$
models and $\epsilon\ge0.1$, against 25\% for SCP+MPPI on the same layouts (Table~\ref{tab:pooled}: $-29.8$ points,
$p=2\times10^{-43}$), at the same return. Pure SCP does not fail for want of convergence: on PointGoal1's $c=1.5$ model
it certifies nearly as many steps as SCP+MPPI (62--95\% against 63--97\% across $\epsilon$; Sec.~\ref{sec:certificates}). It fails
as a local method does, converging to the local solution nearest a poor start: it converges with barrier slack in use,
meaning no safe plan nearby, on more steps (2--19\% against 0--11\% on PointGoal1) and finds a plan that meets the
barrier in the model less often (88\% of steps against 92\%). On the $c=1.5$ models, SCP+MPPI starts from MPPI's plan on
33--43\% of steps on PointGoal1 and 38--55\% on PointGoal2 ($\epsilon\ge0.1$).

\textbf{With a sampled start, SCP+MPPI is as safe as sampling and earns more than its proposal.} Against MPPI and CEM
with the 512 sequences that seed it, SCP+MPPI's violation rate differs by $+1.9$ and $+2.6$ points ($p=0.24$ and
$0.098$), while its return is higher by 0.21 and 0.18 with intervals that exclude zero. Against samplers given the same
time, with 128 times as many sequences, the differences in violations ($-0.1$ and $+0.7$ points) and return ($+0.06$
and $+0.09$) are small and not significant. At $N=10$, then, SCP's refinement adds about as much return as 128 times
more samples, and neither makes the planner safer: with a two-dimensional action, sampling finds plans that meet the
barrier about as often as optimisation does (Sec.~\ref{sec:latent_safety}). The per-model results agree
(Table~\ref{tab:per_model}): no time-matched sampler differs significantly from SCP+MPPI in violations on any model,
and the one significant difference against the 512-sequence samplers favours MPPI on PointGoal2 with $c=1.2$ (23\%
against 34\%, Holm-corrected $p=0.031$).

\textbf{The margin is the main safety lever.} Without a margin ($\epsilon=0$) every planner fails in 30--63\% of
episodes, and pure SCP in 72--85\% (Fig.~\ref{fig:eps_sweep}): the barrier holds only in the model, so a margin must
absorb the model's error. A margin of 0.1--0.3\,m cuts SCP+MPPI's violation rate from 40--50\% to 15--23\% at a moderate
cost in return. At $\epsilon=0.5$\,m the planners turn conservative and return collapses (SCP+MPPI on PointGoal1: 0.38
against 1.75 at $\epsilon=0.1$), and on the more cluttered PointGoal2 violations rise again for every planner
(40--45\%; pure SCP 70\%) as the barrier becomes infeasible in many states, which SCP reports (Sec.~\ref{sec:certificates}).

% ---------------------------------------------------------------------------------------------------------------------
\subsection{At longer horizons SCP+MPPI earns more than time-matched sampling}\label{sec:horizon}

\begin{figure}[t]
  \centering
  \includegraphics[width=\linewidth]{\figdir/main_text/fig_horizon.pdf}
  \caption{Violation rate and return against the horizon $N$ on PointGoal1 ($c=1.5$, $\epsilon=0.2$), each sampler at
  its horizon's time-matched budget, with 95\% intervals. Each $N$ uses a world model trained on windows of $N+1$
  steps; 30 episodes per point (40 for SCP+MPPI and SCP at $N=10$).}
  \label{fig:horizon}
\end{figure}

% CUT 1 if over the page limit: move this figure to the appendix and keep the paragraph's numbers.
\begin{figure}[t]
  \centering
  \includegraphics[width=\linewidth]{\figdir/main_text/fig_return_vs_compute.pdf}
  \caption{Return against mean plan time per step at $N=20$ and $30$, with 95\% intervals. Each sampler is shown at
  every budget we ran (sequences per iteration along the bottom); SCP+MPPI and SCP are single points.}
  \label{fig:compute}
\end{figure}

We repeat the comparison on PointGoal1 ($c=1.5$, $\epsilon=0.2$) at $N=20$ and $30$, re-matching the samplers' budgets
in time at each horizon (Fig.~\ref{fig:horizon}). At both horizons SCP+MPPI earns more than the time-matched samplers,
at the same safety (2 and 1 violations in 30 episodes, against 3 and 0). Its return is higher by 0.65 [0.03, 1.30]
(MPPI, $p=0.021$) and 0.93 [0.15, 1.69] (CEM, $p=0.009$) at $N=20$, and by 1.13 [0.41, 1.93] ($p<0.001$) and 0.98
[0.08, 1.93] ($p=0.005$) at $N=30$.

The samplers do not turn extra time into return (Fig.~\ref{fig:compute}). Their largest budgets earn less than smaller
ones at both horizons, and their best budget changes with the sampler and the horizon (512 sequences for CEM at
$N=20$, 32,768 at $N=30$). At its best budget each sampler ties with SCP+MPPI (CEM: $-0.24$ at $N=20$, $+0.11$ at
$N=30$; neither significant), but that budget is not known in advance. SCP+MPPI used the same settings at every
horizon, and its refinement adds 0.27 (not significant) and 0.67 ($p=0.043$) to the return of the 512-sequence MPPI
plan it starts from.
% Possible reason for the samplers' trend, not tested: the elite set grows with J (top-k = J/8), so a larger budget
% refits to the same top eighth of the samples instead of concentrating on the best of them.

Two caveats apply. The horizon study uses one world model per horizon, one margin and 30 seeds, and the model changes
with $N$, so returns compare within a horizon but not across horizons. And within its 40-iteration cap SCP+MPPI
converged on 53\% of steps at $N=20$ and 32\% at $N=30$, executing its last accepted iterate on the rest: the
comparison is at matched time, not at convergence.

% ---------------------------------------------------------------------------------------------------------------------
\subsection{SCP's convergence certifies the plan and detects infeasibility}\label{sec:certificates}

\begin{figure}[t]
  \centering
  \includegraphics[width=\linewidth]{\figdir/main_text/fig_scp_stop_reasons.pdf}
  \caption{How SCP+MPPI's steps ended on each task's $c=1.5$ model at $N=10$: converged with the plan's rollout meeting
  the barrier (certified), converged with barrier slack in use, or at the 40-iteration cap.}
  \label{fig:stops}
\end{figure}

Each SCP step ends in one of three ways (Fig.~\ref{fig:stops}). Because SCP's reference trajectory is always the model
rollout of its plan, a step that stops as converged, with prox-gradient norm below $10^{-2}$ and barrier excess below
$10^{-4}$, executes an approximately stationary plan that meets the barrier at every step in the model: convergence
itself certifies the plan, at no extra cost. This is the most common outcome. On the $c=1.5$ models SCP+MPPI certifies
97\% of steps at $\epsilon=0$, falling to 63\% at $\epsilon=0.5$, on PointGoal1, and 83\% falling to 60\% on
PointGoal2, using 9--22 iterations on average.

A step that converges with barrier slack in use signals something different: a stationary point at which the barrier
cannot be met, that is, no safe plan near the current one. Its share grows with the margin, from 0 to 11\% of steps on
PointGoal1 and from 1\% to 22\% on PointGoal2, where at $\epsilon=0.5$ it coincides with the rise in violations
(Sec.~\ref{sec:q1}). Sampling planners give no such signal: when no sampled sequence meets the barrier they cannot tell
an empty safe set from too few samples. The remaining steps reach the iteration cap, a share that grows with $\epsilon$
(3\% to 27\% on PointGoal1) and with the horizon (46\% at $N=20$ and 68\% at $N=30$ on the same task and model bound).

% ---------------------------------------------------------------------------------------------------------------------
\subsection{Safety in the model predicts safety in the environment}\label{sec:latent_safety}

% CUT 2 if over the page limit: move this figure to the appendix; the paragraph stands on its numbers.
\begin{figure}[t]
  \centering
  \includegraphics[width=\linewidth]{\figdir/main_text/fig_latent_safety.pdf}
  \caption{Share of steps followed by a violation within $N$ steps, after a plan whose model rollout met the barrier
  and after one whose rollout did not, over the six $N=10$ models and $\epsilon\ge0.1$, with 95\% Wilson intervals
  (log scale).}
  \label{fig:latent_safety}
\end{figure}

The barrier is enforced only in the learned model; Fig.~\ref{fig:latent_safety} tests whether that is enough. After a
plan that meets the barrier in the model, a violation follows within $N$ steps on $2$--$10\times10^{-4}$ of steps,
whichever planner made the plan (including the 512-sequence samplers); after a plan that does not, on
$3$--$7\times10^{-2}$, 56 to 131 times as often. Conversely, 93--98\% of violations were preceded within $N$ steps by a
plan that broke the barrier in the model: the model nearly always warns before a violation. Violations therefore arise
mainly in states where the planner finds no plan that is safe in the model, not from plans that are safe in the model
but not in the environment. All planners find such a plan on a similar share of steps (90--92\%; pure SCP 88\%), which
is why optimising harder does not make the planner safer at $N=10$ (Sec.~\ref{sec:q1}).

% ---------------------------------------------------------------------------------------------------------------------
\subsection{Regularity constraints}\label{sec:regularity_results}

Bounding the networks' Lipschitz constants at $c=1.5$ costs no measurable safety or return: against the unbounded
models, no planner's violation rate or return differs significantly on either task (McNemar $p\ge0.18$; return
intervals include zero;
Figs.~\ref{fig:grid_violations} and~\ref{fig:grid_return} in the appendix). It does make SCP converge more reliably
(Fig.~\ref{fig:lipschitz_convergence}). At $\epsilon$ from 0.1 to 0.3, SCP+MPPI reaches the iteration cap on 6--16\% of
steps on PointGoal1 and 16--21\% on PointGoal2 with $c=1.5$, against 33--44\% and 31--35\% unbounded, needs fewer
iterations (12--16 against 24--28 on PointGoal1), and so certifies more of its plans (78--92\% against 55--65\% on
PointGoal1; 74--83\% against 53--62\% on PointGoal2). The tighter bound $c=1.2$ makes the model less conservative:
returns are higher, but so are violations, significantly on PointGoal1 for SCP+MPPI (44 against 28 of 160 episodes,
$p=0.044$), MPPI ($p=0.009$) and CEM ($p=0.005$), and on PointGoal2 for SCP+MPPI (54 against 38, $p=0.036$). We
therefore use $c=1.5$ for the main results.
. Appendix: % =====================================================================================================================
% Appendix results: \input{appendix_results} in the appendix. Files and numbers as in experiments_results.tex.
% Order is by importance; drop subsections from the end if the appendix must be short.
% =====================================================================================================================
\providecommand{\figdir}{../paper_graphs}

\section{Additional Results}\label{app:results}

% ---------------------------------------------------------------------------------------------------------------------
\subsection{Per-model results}\label{app:per_model}

\begin{table}[htbp]
  \centering
  \caption{Table~\ref{tab:pooled} on the $c=1.2$ world models.}
  \label{tab:pooled_c1_2}
  \small
  \input{\figdir/appendix/table_pooled_c1_2}
\end{table}

\begin{table}[htbp]
  \centering
  \caption{Table~\ref{tab:pooled} on the unbounded ($c=\infty$) world models.}
  \label{tab:pooled_cinf}
  \small
  \input{\figdir/appendix/table_pooled_cinf}
\end{table}

Table~\ref{tab:per_model} gives each planner's violation rate and return on every $N=10$ world model, pooled over
$\epsilon\ge0.1$, and Figs.~\ref{fig:grid_violations} and~\ref{fig:grid_return} show the full sweeps. Pure SCP is
significantly less safe than SCP+MPPI on every model. Against the samplers, the only significant difference is MPPI
with 512 sequences on PointGoal2 with $c=1.2$, which violates less.

\begin{table}[htbp]
  \centering
  \caption{Each planner on each $N=10$ world model, pooled over $\epsilon\ge0.1$: violation rate (Wilson) and return
  (bootstrap) with 95\% intervals, and McNemar's $p$ for SCP+MPPI against it on paired seeds, Holm-corrected over the
  six models.}
  \label{tab:per_model}
  \small
  \input{\figdir/appendix/table_per_model}
\end{table}

\begin{figure}[htbp]
  \centering
  \includegraphics[width=\linewidth]{\figdir/appendix/fig_eps_sweep_all_bounds_violations.pdf}
  \caption{Violation rate against $\epsilon$ for every planner on every $N=10$ world model (rows: tasks; columns:
  Lipschitz bound $c$), with 95\% Wilson intervals.}
  \label{fig:grid_violations}
\end{figure}

\begin{figure}[htbp]
  \centering
  \includegraphics[width=\linewidth]{\figdir/appendix/fig_eps_sweep_all_bounds_return.pdf}
  \caption{Return against $\epsilon$ for every planner on every $N=10$ world model, with 95\% bootstrap intervals.}
  \label{fig:grid_return}
\end{figure}

% ---------------------------------------------------------------------------------------------------------------------
\subsection{Lipschitz bound and SCP's convergence}\label{app:lipschitz}

\begin{figure}[htbp]
  \centering
  \includegraphics[width=\linewidth]{\figdir/appendix/fig_lipschitz_convergence.pdf}
  \caption{Share of SCP+MPPI's steps that reached the 40-iteration cap against $\epsilon$, for each Lipschitz bound,
  at $N=10$.}
  \label{fig:lipschitz_convergence}
\end{figure}

Without a Lipschitz bound, SCP+MPPI reaches the iteration cap on more steps than with either bound at every margin,
except at $\epsilon=0.5$ on PointGoal2, where $c=1.5$ and $c=\infty$ are within 0.02 (Fig.~\ref{fig:lipschitz_convergence}).
Which bound converges best differs by task: $c=1.5$ on PointGoal1, $c=1.2$ on PointGoal2.

% ---------------------------------------------------------------------------------------------------------------------
\subsection{What more samples buy at \texorpdfstring{$N=10$}{N=10}}\label{app:samples}

\begin{figure}[htbp]
  \centering
  \includegraphics[width=\linewidth]{\figdir/ablation/fig_sample_vs_time_matched.pdf}
  \caption{Fig.~\ref{fig:eps_sweep} with the samplers at both budgets: 512 sequences (open markers, dashed) and
  time-matched (filled, solid).}
  \label{fig:sample_vs_time_matched}
\end{figure}

\begin{figure}[htbp]
  \centering
  \includegraphics[width=\linewidth]{\figdir/ablation/fig_return_vs_violations.pdf}
  \caption{Return against violation rate on each task's $c=1.5$ model at $N=10$, pooled over $\epsilon\ge0.1$, with
  95\% intervals. Arrows lead from each sampler with 512 sequences to the same sampler time-matched.}
  \label{fig:trade_off}
\end{figure}

Figs.~\ref{fig:sample_vs_time_matched} and~\ref{fig:trade_off} show what 128 times as many sequences buy the samplers.
On PointGoal1's $c=1.5$ model, the time-matched samplers earn more than with 512 sequences but violate more often
(MPPI: return 0.78 to 1.15, violation rate 0.19 to 0.24; CEM: 0.84 to 0.94 and 0.14 to 0.26), while SCP+MPPI reaches
the higher return at the lower violation rate (1.15 and 0.17). On PointGoal2 the extra sequences raise return (0.35 to
0.51 and 0.49) without changing the violation rate. These shifts lie within the intervals; the paired tests of
Sec.~\ref{sec:q1} find no significant difference in safety between SCP+MPPI and either budget.

% ---------------------------------------------------------------------------------------------------------------------
\subsection{All budgets at every horizon}\label{app:horizon_budgets}

\begin{figure}[htbp]
  \centering
  \includegraphics[width=\linewidth]{\figdir/ablation/fig_horizon_all_budgets.pdf}
  \caption{Fig.~\ref{fig:horizon} with the samplers at both 512 sequences (open markers, dashed) and their
  time-matched budget (filled, solid).}
  \label{fig:horizon_all}
\end{figure}

With 512 sequences the samplers' return at $N=20$ is close to SCP+MPPI's (MPPI 2.07, CEM 2.58, against 2.34; neither
difference significant), and at $N=30$ below it (1.59 and 1.22 against 2.26; $p=0.043$ and $0.013$)
(Fig.~\ref{fig:horizon_all}). Together with Fig.~\ref{fig:compute}, this shows that the samplers' best budget depends on
the sampler and the horizon.

% ---------------------------------------------------------------------------------------------------------------------
\subsection{Does optimisation exploit the model?}\label{app:optimism}

\begin{figure}[htbp]
  \centering
  \includegraphics[width=\linewidth]{\figdir/ablation/fig_model_optimism.pdf}
  \caption{Model optimism: each executed plan's predicted return (discounted predicted rewards plus terminal value)
  minus the discounted return received from that step on, averaged over each episode, with 95\% bootstrap intervals.
  $N=10$ pools every world model ($\epsilon\ge0.1$); $N=20$ and $30$ are the horizon study's models.}
  \label{fig:optimism}
\end{figure}

A gradient-based optimiser might exploit model error, finding plans whose predicted return the environment does not
deliver. We measure each executed plan's optimism: its predicted return, whose terminal value carries the goal bonus,
minus the discounted return actually received from that step on, using steps at least 299 steps before the episode
ends so that the part of the future the episode did not reach weighs under 5\%. Every planner is optimistic by about
2, as expected of a value fitted to the 0.7 expectile, but by nearly the same amount (Fig.~\ref{fig:optimism}): 2.06
(pure SCP) to 2.19 at $N=10$, 2.20 to 2.44 at $N=20$ and 2.41 to 2.56 at $N=30$, with SCP+MPPI inside the samplers'
range at every horizon and no trend with the sampling budget. We find no evidence that SCP, or a larger sampling
budget, exploits the model more than sampling with 512 sequences does.

% ---------------------------------------------------------------------------------------------------------------------
\subsection{Warned and unwarned violations}\label{app:crashes}

\begin{figure}[htbp]
  \centering
  \includegraphics[width=\linewidth]{\figdir/ablation/fig_crash_types.pdf}
  \caption{Violations per 100 episodes over the six $N=10$ models and $\epsilon\ge0.1$, split into those preceded
  within $N$ steps by a plan that broke the barrier in the model (warned, light) and those that were not (unwarned,
  dark, counts above the bars).}
  \label{fig:crash_types}
\end{figure}

Few violations come without warning (Fig.~\ref{fig:crash_types}). SCP+MPPI has 14 unwarned violations in 880
episodes, against 6 and 7 for the 512-sequence samplers and 7 and 13 for the time-matched ones. That stronger
optimisation, by SCP or by more samples, leads to more violations the model did not foresee would be consistent with
it planning closer to where the model is wrong, but paired by seed none of these differences is significant
($p\ge0.12$).

% ---------------------------------------------------------------------------------------------------------------------
\subsection{When violations happen}\label{app:survival}

\begin{figure}[htbp]
  \centering
  \includegraphics[width=\linewidth]{\figdir/appendix/fig_survival_pg1_c1_5.pdf}
  \caption{Share of episodes without a violation by each step, at each margin, on PointGoal1's $c=1.5$ model at
  $N=10$. Episodes end early only at a violation, so these are Kaplan--Meier curves without censoring.}
  \label{fig:survival}
\end{figure}

Most violations happen within the first 300 or so steps, after which the curves flatten (Fig.~\ref{fig:survival});
pure SCP's keeps falling for longer.
. Citations: experiments_refs.bib.
%
% Every figure and table is written by the repository's scripts, run from the repository root:
%   python src/scripts/analyse_results.py  -> src/save_data/experiments/analysis_results/*.csv
%   python src/scripts/paper_figures.py    -> paper_graphs/{main_text,ablation,appendix}/ (figures and both tables)
%   python src/scripts/env_figure.py       -> paper_graphs/main_text/fig_environments.pdf
% \figdir is the paper_graphs folder; define it before \input if the paper lives elsewhere (e.g. on Overleaf).
%
% Notation follows the method: N horizon, \epsilon barrier margin, \alpha DCBF rate, c Lipschitz bound, \tau risk
% threshold, g_\Phi margin head. Labels of the method's parts used here (rename to match yours): sec:regularity,
% sec:cbf, sec:scp, sec:init, eq:ocp.
% Every number is from the analysis of the runs on main (2d4240b). "% TODO" marks a fact to confirm before submission.
% =====================================================================================================================
\providecommand{\figdir}{../paper_graphs}

\section{Experimental Setup}\label{sec:setup}

Our experiments answer four questions.
\textbf{(Q1)}~Does SCP, started from a sampled plan, plan better than sampling-based MPC given the same wall-clock time?
\textbf{(Q2)}~How does the comparison change as the horizon, and with it the number of decision variables, grows?
\textbf{(Q3)}~What does SCP's convergence tell us online?
\textbf{(Q4)}~Does meeting the barrier in the learned model keep the robot safe in the environment, and what do the
regularity constraints cost?

% ---------------------------------------------------------------------------------------------------------------------
\subsection{Environments}\label{sec:envs}

\begin{figure}[t]
  \centering
  \includegraphics[width=0.62\linewidth]{\figdir/main_text/fig_environments.pdf}
  \caption{The two tasks from above, one layout each (layouts are drawn from the episode's seed). The point robot
  (red) must reach the goal (green) without entering a hazard (blue); in PointGoal2, touching a vase (cyan) is also a
  violation.}
  \label{fig:envs}
\end{figure}

We use \texttt{SafetyPointGoal1} and \texttt{SafetyPointGoal2} from Safety-Gymnasium~\citep{ji2023safety}
(Fig.~\ref{fig:envs}). A point robot with two actuators (forward force and turning, $u\in[-1,1]^2$) must reach a goal
of radius $0.3$\,m, which moves to a new position each time it is reached, within $T=1000$ steps. The reward is the
per-step reduction in distance to the goal plus a bonus of 1 for reaching it. The robot observes 60 quantities: 12
proprioceptive signals (accelerometer, velocimeter, gyroscope and magnetometer) and three 16-ray pseudo-lidars measuring
the proximity of the goal, the hazards and the vases. Level~1 places 8 hazards of radius $0.2$\,m and one cost-free
vase in a $3\times3$\,m region; level~2 places 10 hazards and 10 vases in a $4\times4$\,m region and also charges for
touching or moving a vase. A step is a \emph{violation} if it incurs any cost. We treat violations as absorbing
failures: an episode ends at its first violation, so the \emph{violation rate} is the share of episodes that fail and
the \emph{return} is the reward collected before failure or the time limit.

These tasks suit our question for four reasons.
(i)~DSRL~\citep{liu2023datasets} provides offline datasets for both, so the world model is learned without any
interaction, which is the setting of this paper.
(ii)~The constraint is geometric (stay out of regions of the plane) and observed through the lidar, so the barrier's
margin is a distance in metres, the margin $\epsilon$ has a physical meaning, and safety in the model can be checked
against safety in the environment (Sec.~\ref{sec:latent_safety}).
(iii)~Level~2 is more cluttered (20 obstacles in 16\,m$^2$ against 8 in 9\,m$^2$) and adds a second, movable kind of
obstacle, so the two levels, together with $\epsilon$, vary how restrictive the safe set is.
(iv)~The action is two-dimensional, which favours sampling: at $N=10$ a plan has only 20 decision variables, so the
comparison is a demanding one for SCP. The horizon study (Sec.~\ref{sec:horizon}) raises this to 40 and 60.

% ---------------------------------------------------------------------------------------------------------------------
\subsection{Offline data and world models}\label{sec:wm_setup}

For each task we train on the corresponding DSRL dataset (2,022 trajectories for PointGoal1 and 3,442 for PointGoal2,
of up to 1000 steps each), truncating every trajectory at its first violation to match the absorbing failure above.
This leaves about 0.87\,M transitions for PointGoal1 and 1.15\,M for PointGoal2; only 329 of PointGoal2's trajectories
run their full length without a violation.
% TODO: PG1's exact transition count (FilterDataset.ipynb prints it for PG2 only: 1,145,452); one sentence on how DSRL
% collected the data (behaviour policies, cost budgets).
The encoder drops seven proprioceptive channels (the accelerometer and the out-of-plane components of velocity,
angular rate and magnetic field) and maps the remaining 53 to a latent state of $d_z=32$ dimensions. The four networks
the planner differentiates through (dynamics, reward, value and margin) are MLPs with two hidden layers of 256 units,
each with Lipschitz constant at most $c$ (Sec.~\ref{sec:regularity}). We train one model per task and bound,
$c\in\{1.2,1.5,\infty\}$ ($c=\infty$: unbounded), for $N=10$, and two more PointGoal1 models with $c=1.5$ for $N=20$
and $30$. Each is trained for 250 epochs of 1000 (PointGoal1) or 1500 (PointGoal2) minibatches of 256 windows of $N+1$
consecutive transitions, with learning rate $3\times10^{-4}$, the loss at rollout step $t$ weighted by $\rho^t$
($\rho=0.9$, $0.95$, $0.99$ for $N=10$, $20$, $30$), discount $\gamma=0.99$ and value expectile $0.7$. The margin head
$g_\Phi$ regresses the distance from the robot to the nearest hazard edge (in PointGoal2, hazard or vase edge),
computed from the lidar, clipped at 1\,m, with samples near an edge up-weighted.
% TODO: optimiser (Adam?); state that each model is a single training seed (limitations).

% ---------------------------------------------------------------------------------------------------------------------
\subsection{Planners}\label{sec:planners}

All planners solve the same receding-horizon problem~\eqref{eq:ocp} in the same model: maximise
$\sum_{k=0}^{N-1}\gamma^k r(z_k,u_k)+\gamma^N V(z_N)-\rho_u\|u\|^2$ ($\rho_u=10^{-4}$) subject to
$(1-\alpha)h(z_k)-h(z_{k+1})\le\tau=0$ at every step, with $h(z)=g_\Phi(z)-\epsilon$ and $\alpha=0.05$. Each executes
the first action of its plan and replans at every step.

\textbf{SCP} is the method of Sec.~\ref{sec:scp} started from its previous plan shifted by one step (zeros at the first
step): at most 40 prox-linear iterations, each a convex QP solved by Clarabel~\citep{goulart2024clarabel}, with exact
$\ell_1$ penalties $w_{\mathrm{dyn}}=10^3$ and $w_{\mathrm{safe}}=10^2$, a trust-region weight that starts at 1, is
floored at $10^{-2}$ and is carried between steps, acceptance thresholds $(0.25,0.75)$, and a stop when the
prox-gradient norm falls below $10^{-2}$.
\textbf{SCP+MPPI}, our method, first runs MPPI with 512 sequences (below) and starts SCP from whichever of MPPI's plan
and its own shifted plan has the lower penalty merit (Sec.~\ref{sec:init}). SCP accepts only steps that lower the
merit, so the plan SCP+MPPI executes is never worse under the merit than MPPI's, whether or not SCP converges.
\textbf{MPPI}~\citep{williams2017information} and \textbf{CEM}~\citep{rubinstein1999cross} run 6 iterations of $J$
action sequences drawn from a Gaussian that is refitted to the best $J/8$ (weighted by $\exp(\text{score}/0.5)$ in MPPI,
uniformly in CEM) and warm-started from the previous mean shifted by one step. A sequence whose rollout breaks the
barrier at any step ranks below every sequence that meets it; if none meets it, sequences are ranked by their total
barrier excess, the quantity SCP's $\ell_1$ penalty charges. The executed plan is the refitted mean if its rollout
meets the barrier, and otherwise the best sequence sampled.

We run the samplers with $J=512$, the proposal SCP+MPPI uses, and \emph{time-matched}, with $J$ chosen so that a step
takes about as long as SCP+MPPI's: $J=65{,}536$ at $N=10$ (about 0.70\,s against 0.40--0.71\,s for SCP+MPPI, timed on
the same idle GPU), and $95{,}040$ and $120{,}000$ at $N=20$ and $30$ (mean logged plan times 2.04\,s and 3.81\,s,
against 2.16\,s and 3.61\,s for SCP+MPPI). Each run used one 1g.10gb slice of an NVIDIA A100, with the QP solved on the
CPU.
% TODO: the 65,536 time is interpolated between 32,768 (0.33 s) and 131,072 (1.43 s) sequences; time it directly or
% say so. Confirm the hardware sentence.

% ---------------------------------------------------------------------------------------------------------------------
\subsection{Protocol and statistics}\label{sec:protocol}

We sweep the margin $\epsilon\in\{0,0.1,0.2,0.3,0.5\}$\,m on every $N=10$ model and fix $\epsilon=0.2$\,m for the
horizon study. Each configuration runs on seeds $0$--$39$, or $0$--$29$ for the time-matched samplers, the horizon
study, PointGoal2's $c=\infty$ model and SCP+MPPI on PointGoal1's $c=\infty$ model. The seed fixes the layout, the
initial state and the planner's random numbers, so every comparison between planners or models is paired by seed. We
report violation rates with Wilson 95\% intervals and mean returns with bootstrap 95\% intervals ($10^4$ resamples).
Paired comparisons report the difference in violation rate and in return with a 95\% bootstrap interval, and the
probability, given the data, that SCP+MPPI's violation rate is lower or its return higher than the baseline's, from a
Bayesian bootstrap~\citep{rubin1981bayesian}. When pooling over margins and tasks we use $\epsilon\ge0.1$, since at
$\epsilon=0$ every planner fails often (Sec.~\ref{sec:q1}), and resample whole seeds, since a seed's episodes share a
layout. The per-model results in the appendix use McNemar's exact test on violations, Holm-corrected over the six models,
and the Wilcoxon signed-rank test on returns.

% =====================================================================================================================
\section{Results}\label{sec:results}

% ---------------------------------------------------------------------------------------------------------------------
\subsection{SCP needs a sampled start, and then matches time-matched sampling at \texorpdfstring{$N=10$}{N=10}}\label{sec:q1}

\begin{figure}[t]
  \centering
  \includegraphics[width=\linewidth]{\figdir/main_text/fig_eps_sweep_c1_5.pdf}
  \caption{Violation rate (top) and return (bottom) against the margin $\epsilon$ on each task's $c=1.5$ model at
  $N=10$, with 95\% intervals; 40 episodes per point (30 for the time-matched samplers).}
  \label{fig:eps_sweep}
\end{figure}

\begin{table}[t]
  \centering
  \caption{SCP+MPPI minus each baseline on the $c=1.5$ world models at $N=10$, pooled over both tasks and
  $\epsilon\ge0.1$ and paired by task, margin and seed (Tables~\ref{tab:pooled_c1_2} and~\ref{tab:pooled_cinf} in the
  appendix give the other bounds). A negative $\Delta$ violation rate or a positive $\Delta$ return favours SCP+MPPI.
  Brackets: 95\% intervals resampling whole seeds. $P$(safer) and $P$(higher return): the probability, given the data,
  that SCP+MPPI's violation rate is lower, or its return higher, than the baseline's on these models (Bayesian bootstrap
  over seeds).}
  \label{tab:pooled}
  \small
  \input{\figdir/main_text/table_pooled_c1_5}
\end{table}

\textbf{SCP needs a sampled start.} Started from its own shifted plan, SCP fails in 55\% of episodes over all $N=10$
models and $\epsilon\ge0.1$, against 25\% for SCP+MPPI on the same layouts (Table~\ref{tab:pooled}: $-29.8$ points,
$p=2\times10^{-43}$), at the same return. Pure SCP does not fail for want of convergence: on PointGoal1's $c=1.5$ model
it certifies nearly as many steps as SCP+MPPI (62--95\% against 63--97\% across $\epsilon$; Sec.~\ref{sec:certificates}). It fails
as a local method does, converging to the local solution nearest a poor start: it converges with barrier slack in use,
meaning no safe plan nearby, on more steps (2--19\% against 0--11\% on PointGoal1) and finds a plan that meets the
barrier in the model less often (88\% of steps against 92\%). On the $c=1.5$ models, SCP+MPPI starts from MPPI's plan on
33--43\% of steps on PointGoal1 and 38--55\% on PointGoal2 ($\epsilon\ge0.1$).

\textbf{With a sampled start, SCP+MPPI is as safe as sampling and earns more than its proposal.} Against MPPI and CEM
with the 512 sequences that seed it, SCP+MPPI's violation rate differs by $+1.9$ and $+2.6$ points ($p=0.24$ and
$0.098$), while its return is higher by 0.21 and 0.18 with intervals that exclude zero. Against samplers given the same
time, with 128 times as many sequences, the differences in violations ($-0.1$ and $+0.7$ points) and return ($+0.06$
and $+0.09$) are small and not significant. At $N=10$, then, SCP's refinement adds about as much return as 128 times
more samples, and neither makes the planner safer: with a two-dimensional action, sampling finds plans that meet the
barrier about as often as optimisation does (Sec.~\ref{sec:latent_safety}). The per-model results agree
(Table~\ref{tab:per_model}): no time-matched sampler differs significantly from SCP+MPPI in violations on any model,
and the one significant difference against the 512-sequence samplers favours MPPI on PointGoal2 with $c=1.2$ (23\%
against 34\%, Holm-corrected $p=0.031$).

\textbf{The margin is the main safety lever.} Without a margin ($\epsilon=0$) every planner fails in 30--63\% of
episodes, and pure SCP in 72--85\% (Fig.~\ref{fig:eps_sweep}): the barrier holds only in the model, so a margin must
absorb the model's error. A margin of 0.1--0.3\,m cuts SCP+MPPI's violation rate from 40--50\% to 15--23\% at a moderate
cost in return. At $\epsilon=0.5$\,m the planners turn conservative and return collapses (SCP+MPPI on PointGoal1: 0.38
against 1.75 at $\epsilon=0.1$), and on the more cluttered PointGoal2 violations rise again for every planner
(40--45\%; pure SCP 70\%) as the barrier becomes infeasible in many states, which SCP reports (Sec.~\ref{sec:certificates}).

% ---------------------------------------------------------------------------------------------------------------------
\subsection{At longer horizons SCP+MPPI earns more than time-matched sampling}\label{sec:horizon}

\begin{figure}[t]
  \centering
  \includegraphics[width=\linewidth]{\figdir/main_text/fig_horizon.pdf}
  \caption{Violation rate and return against the horizon $N$ on PointGoal1 ($c=1.5$, $\epsilon=0.2$), each sampler at
  its horizon's time-matched budget, with 95\% intervals. Each $N$ uses a world model trained on windows of $N+1$
  steps; 30 episodes per point (40 for SCP+MPPI and SCP at $N=10$).}
  \label{fig:horizon}
\end{figure}

% CUT 1 if over the page limit: move this figure to the appendix and keep the paragraph's numbers.
\begin{figure}[t]
  \centering
  \includegraphics[width=\linewidth]{\figdir/main_text/fig_return_vs_compute.pdf}
  \caption{Return against mean plan time per step at $N=20$ and $30$, with 95\% intervals. Each sampler is shown at
  every budget we ran (sequences per iteration along the bottom); SCP+MPPI and SCP are single points.}
  \label{fig:compute}
\end{figure}

We repeat the comparison on PointGoal1 ($c=1.5$, $\epsilon=0.2$) at $N=20$ and $30$, re-matching the samplers' budgets
in time at each horizon (Fig.~\ref{fig:horizon}). At both horizons SCP+MPPI earns more than the time-matched samplers,
at the same safety (2 and 1 violations in 30 episodes, against 3 and 0). Its return is higher by 0.65 [0.03, 1.30]
(MPPI, $p=0.021$) and 0.93 [0.15, 1.69] (CEM, $p=0.009$) at $N=20$, and by 1.13 [0.41, 1.93] ($p<0.001$) and 0.98
[0.08, 1.93] ($p=0.005$) at $N=30$.

The samplers do not turn extra time into return (Fig.~\ref{fig:compute}). Their largest budgets earn less than smaller
ones at both horizons, and their best budget changes with the sampler and the horizon (512 sequences for CEM at
$N=20$, 32,768 at $N=30$). At its best budget each sampler ties with SCP+MPPI (CEM: $-0.24$ at $N=20$, $+0.11$ at
$N=30$; neither significant), but that budget is not known in advance. SCP+MPPI used the same settings at every
horizon, and its refinement adds 0.27 (not significant) and 0.67 ($p=0.043$) to the return of the 512-sequence MPPI
plan it starts from.
% Possible reason for the samplers' trend, not tested: the elite set grows with J (top-k = J/8), so a larger budget
% refits to the same top eighth of the samples instead of concentrating on the best of them.

Two caveats apply. The horizon study uses one world model per horizon, one margin and 30 seeds, and the model changes
with $N$, so returns compare within a horizon but not across horizons. And within its 40-iteration cap SCP+MPPI
converged on 53\% of steps at $N=20$ and 32\% at $N=30$, executing its last accepted iterate on the rest: the
comparison is at matched time, not at convergence.

% ---------------------------------------------------------------------------------------------------------------------
\subsection{SCP's convergence certifies the plan and detects infeasibility}\label{sec:certificates}

\begin{figure}[t]
  \centering
  \includegraphics[width=\linewidth]{\figdir/main_text/fig_scp_stop_reasons.pdf}
  \caption{How SCP+MPPI's steps ended on each task's $c=1.5$ model at $N=10$: converged with the plan's rollout meeting
  the barrier (certified), converged with barrier slack in use, or at the 40-iteration cap.}
  \label{fig:stops}
\end{figure}

Each SCP step ends in one of three ways (Fig.~\ref{fig:stops}). Because SCP's reference trajectory is always the model
rollout of its plan, a step that stops as converged, with prox-gradient norm below $10^{-2}$ and barrier excess below
$10^{-4}$, executes an approximately stationary plan that meets the barrier at every step in the model: convergence
itself certifies the plan, at no extra cost. This is the most common outcome. On the $c=1.5$ models SCP+MPPI certifies
97\% of steps at $\epsilon=0$, falling to 63\% at $\epsilon=0.5$, on PointGoal1, and 83\% falling to 60\% on
PointGoal2, using 9--22 iterations on average.

A step that converges with barrier slack in use signals something different: a stationary point at which the barrier
cannot be met, that is, no safe plan near the current one. Its share grows with the margin, from 0 to 11\% of steps on
PointGoal1 and from 1\% to 22\% on PointGoal2, where at $\epsilon=0.5$ it coincides with the rise in violations
(Sec.~\ref{sec:q1}). Sampling planners give no such signal: when no sampled sequence meets the barrier they cannot tell
an empty safe set from too few samples. The remaining steps reach the iteration cap, a share that grows with $\epsilon$
(3\% to 27\% on PointGoal1) and with the horizon (46\% at $N=20$ and 68\% at $N=30$ on the same task and model bound).

% ---------------------------------------------------------------------------------------------------------------------
\subsection{Safety in the model predicts safety in the environment}\label{sec:latent_safety}

% CUT 2 if over the page limit: move this figure to the appendix; the paragraph stands on its numbers.
\begin{figure}[t]
  \centering
  \includegraphics[width=\linewidth]{\figdir/main_text/fig_latent_safety.pdf}
  \caption{Share of steps followed by a violation within $N$ steps, after a plan whose model rollout met the barrier
  and after one whose rollout did not, over the six $N=10$ models and $\epsilon\ge0.1$, with 95\% Wilson intervals
  (log scale).}
  \label{fig:latent_safety}
\end{figure}

The barrier is enforced only in the learned model; Fig.~\ref{fig:latent_safety} tests whether that is enough. After a
plan that meets the barrier in the model, a violation follows within $N$ steps on $2$--$10\times10^{-4}$ of steps,
whichever planner made the plan (including the 512-sequence samplers); after a plan that does not, on
$3$--$7\times10^{-2}$, 56 to 131 times as often. Conversely, 93--98\% of violations were preceded within $N$ steps by a
plan that broke the barrier in the model: the model nearly always warns before a violation. Violations therefore arise
mainly in states where the planner finds no plan that is safe in the model, not from plans that are safe in the model
but not in the environment. All planners find such a plan on a similar share of steps (90--92\%; pure SCP 88\%), which
is why optimising harder does not make the planner safer at $N=10$ (Sec.~\ref{sec:q1}).

% ---------------------------------------------------------------------------------------------------------------------
\subsection{Regularity constraints}\label{sec:regularity_results}

Bounding the networks' Lipschitz constants at $c=1.5$ costs no measurable safety or return: against the unbounded
models, no planner's violation rate or return differs significantly on either task (McNemar $p\ge0.18$; return
intervals include zero;
Figs.~\ref{fig:grid_violations} and~\ref{fig:grid_return} in the appendix). It does make SCP converge more reliably
(Fig.~\ref{fig:lipschitz_convergence}). At $\epsilon$ from 0.1 to 0.3, SCP+MPPI reaches the iteration cap on 6--16\% of
steps on PointGoal1 and 16--21\% on PointGoal2 with $c=1.5$, against 33--44\% and 31--35\% unbounded, needs fewer
iterations (12--16 against 24--28 on PointGoal1), and so certifies more of its plans (78--92\% against 55--65\% on
PointGoal1; 74--83\% against 53--62\% on PointGoal2). The tighter bound $c=1.2$ makes the model less conservative:
returns are higher, but so are violations, significantly on PointGoal1 for SCP+MPPI (44 against 28 of 160 episodes,
$p=0.044$), MPPI ($p=0.009$) and CEM ($p=0.005$), and on PointGoal2 for SCP+MPPI (54 against 38, $p=0.036$). We
therefore use $c=1.5$ for the main results.
. Appendix: % =====================================================================================================================
% Appendix results: % =====================================================================================================================
% Appendix results: \input{appendix_results} in the appendix. Files and numbers as in experiments_results.tex.
% Order is by importance; drop subsections from the end if the appendix must be short.
% =====================================================================================================================
\providecommand{\figdir}{../paper_graphs}

\section{Additional Results}\label{app:results}

% ---------------------------------------------------------------------------------------------------------------------
\subsection{Per-model results}\label{app:per_model}

\begin{table}[htbp]
  \centering
  \caption{Table~\ref{tab:pooled} on the $c=1.2$ world models.}
  \label{tab:pooled_c1_2}
  \small
  \input{\figdir/appendix/table_pooled_c1_2}
\end{table}

\begin{table}[htbp]
  \centering
  \caption{Table~\ref{tab:pooled} on the unbounded ($c=\infty$) world models.}
  \label{tab:pooled_cinf}
  \small
  \input{\figdir/appendix/table_pooled_cinf}
\end{table}

Table~\ref{tab:per_model} gives each planner's violation rate and return on every $N=10$ world model, pooled over
$\epsilon\ge0.1$, and Figs.~\ref{fig:grid_violations} and~\ref{fig:grid_return} show the full sweeps. Pure SCP is
significantly less safe than SCP+MPPI on every model. Against the samplers, the only significant difference is MPPI
with 512 sequences on PointGoal2 with $c=1.2$, which violates less.

\begin{table}[htbp]
  \centering
  \caption{Each planner on each $N=10$ world model, pooled over $\epsilon\ge0.1$: violation rate (Wilson) and return
  (bootstrap) with 95\% intervals, and McNemar's $p$ for SCP+MPPI against it on paired seeds, Holm-corrected over the
  six models.}
  \label{tab:per_model}
  \small
  \input{\figdir/appendix/table_per_model}
\end{table}

\begin{figure}[htbp]
  \centering
  \includegraphics[width=\linewidth]{\figdir/appendix/fig_eps_sweep_all_bounds_violations.pdf}
  \caption{Violation rate against $\epsilon$ for every planner on every $N=10$ world model (rows: tasks; columns:
  Lipschitz bound $c$), with 95\% Wilson intervals.}
  \label{fig:grid_violations}
\end{figure}

\begin{figure}[htbp]
  \centering
  \includegraphics[width=\linewidth]{\figdir/appendix/fig_eps_sweep_all_bounds_return.pdf}
  \caption{Return against $\epsilon$ for every planner on every $N=10$ world model, with 95\% bootstrap intervals.}
  \label{fig:grid_return}
\end{figure}

% ---------------------------------------------------------------------------------------------------------------------
\subsection{Lipschitz bound and SCP's convergence}\label{app:lipschitz}

\begin{figure}[htbp]
  \centering
  \includegraphics[width=\linewidth]{\figdir/appendix/fig_lipschitz_convergence.pdf}
  \caption{Share of SCP+MPPI's steps that reached the 40-iteration cap against $\epsilon$, for each Lipschitz bound,
  at $N=10$.}
  \label{fig:lipschitz_convergence}
\end{figure}

Without a Lipschitz bound, SCP+MPPI reaches the iteration cap on more steps than with either bound at every margin,
except at $\epsilon=0.5$ on PointGoal2, where $c=1.5$ and $c=\infty$ are within 0.02 (Fig.~\ref{fig:lipschitz_convergence}).
Which bound converges best differs by task: $c=1.5$ on PointGoal1, $c=1.2$ on PointGoal2.

% ---------------------------------------------------------------------------------------------------------------------
\subsection{What more samples buy at \texorpdfstring{$N=10$}{N=10}}\label{app:samples}

\begin{figure}[htbp]
  \centering
  \includegraphics[width=\linewidth]{\figdir/ablation/fig_sample_vs_time_matched.pdf}
  \caption{Fig.~\ref{fig:eps_sweep} with the samplers at both budgets: 512 sequences (open markers, dashed) and
  time-matched (filled, solid).}
  \label{fig:sample_vs_time_matched}
\end{figure}

\begin{figure}[htbp]
  \centering
  \includegraphics[width=\linewidth]{\figdir/ablation/fig_return_vs_violations.pdf}
  \caption{Return against violation rate on each task's $c=1.5$ model at $N=10$, pooled over $\epsilon\ge0.1$, with
  95\% intervals. Arrows lead from each sampler with 512 sequences to the same sampler time-matched.}
  \label{fig:trade_off}
\end{figure}

Figs.~\ref{fig:sample_vs_time_matched} and~\ref{fig:trade_off} show what 128 times as many sequences buy the samplers.
On PointGoal1's $c=1.5$ model, the time-matched samplers earn more than with 512 sequences but violate more often
(MPPI: return 0.78 to 1.15, violation rate 0.19 to 0.24; CEM: 0.84 to 0.94 and 0.14 to 0.26), while SCP+MPPI reaches
the higher return at the lower violation rate (1.15 and 0.17). On PointGoal2 the extra sequences raise return (0.35 to
0.51 and 0.49) without changing the violation rate. These shifts lie within the intervals; the paired tests of
Sec.~\ref{sec:q1} find no significant difference in safety between SCP+MPPI and either budget.

% ---------------------------------------------------------------------------------------------------------------------
\subsection{All budgets at every horizon}\label{app:horizon_budgets}

\begin{figure}[htbp]
  \centering
  \includegraphics[width=\linewidth]{\figdir/ablation/fig_horizon_all_budgets.pdf}
  \caption{Fig.~\ref{fig:horizon} with the samplers at both 512 sequences (open markers, dashed) and their
  time-matched budget (filled, solid).}
  \label{fig:horizon_all}
\end{figure}

With 512 sequences the samplers' return at $N=20$ is close to SCP+MPPI's (MPPI 2.07, CEM 2.58, against 2.34; neither
difference significant), and at $N=30$ below it (1.59 and 1.22 against 2.26; $p=0.043$ and $0.013$)
(Fig.~\ref{fig:horizon_all}). Together with Fig.~\ref{fig:compute}, this shows that the samplers' best budget depends on
the sampler and the horizon.

% ---------------------------------------------------------------------------------------------------------------------
\subsection{Does optimisation exploit the model?}\label{app:optimism}

\begin{figure}[htbp]
  \centering
  \includegraphics[width=\linewidth]{\figdir/ablation/fig_model_optimism.pdf}
  \caption{Model optimism: each executed plan's predicted return (discounted predicted rewards plus terminal value)
  minus the discounted return received from that step on, averaged over each episode, with 95\% bootstrap intervals.
  $N=10$ pools every world model ($\epsilon\ge0.1$); $N=20$ and $30$ are the horizon study's models.}
  \label{fig:optimism}
\end{figure}

A gradient-based optimiser might exploit model error, finding plans whose predicted return the environment does not
deliver. We measure each executed plan's optimism: its predicted return, whose terminal value carries the goal bonus,
minus the discounted return actually received from that step on, using steps at least 299 steps before the episode
ends so that the part of the future the episode did not reach weighs under 5\%. Every planner is optimistic by about
2, as expected of a value fitted to the 0.7 expectile, but by nearly the same amount (Fig.~\ref{fig:optimism}): 2.06
(pure SCP) to 2.19 at $N=10$, 2.20 to 2.44 at $N=20$ and 2.41 to 2.56 at $N=30$, with SCP+MPPI inside the samplers'
range at every horizon and no trend with the sampling budget. We find no evidence that SCP, or a larger sampling
budget, exploits the model more than sampling with 512 sequences does.

% ---------------------------------------------------------------------------------------------------------------------
\subsection{Warned and unwarned violations}\label{app:crashes}

\begin{figure}[htbp]
  \centering
  \includegraphics[width=\linewidth]{\figdir/ablation/fig_crash_types.pdf}
  \caption{Violations per 100 episodes over the six $N=10$ models and $\epsilon\ge0.1$, split into those preceded
  within $N$ steps by a plan that broke the barrier in the model (warned, light) and those that were not (unwarned,
  dark, counts above the bars).}
  \label{fig:crash_types}
\end{figure}

Few violations come without warning (Fig.~\ref{fig:crash_types}). SCP+MPPI has 14 unwarned violations in 880
episodes, against 6 and 7 for the 512-sequence samplers and 7 and 13 for the time-matched ones. That stronger
optimisation, by SCP or by more samples, leads to more violations the model did not foresee would be consistent with
it planning closer to where the model is wrong, but paired by seed none of these differences is significant
($p\ge0.12$).

% ---------------------------------------------------------------------------------------------------------------------
\subsection{When violations happen}\label{app:survival}

\begin{figure}[htbp]
  \centering
  \includegraphics[width=\linewidth]{\figdir/appendix/fig_survival_pg1_c1_5.pdf}
  \caption{Share of episodes without a violation by each step, at each margin, on PointGoal1's $c=1.5$ model at
  $N=10$. Episodes end early only at a violation, so these are Kaplan--Meier curves without censoring.}
  \label{fig:survival}
\end{figure}

Most violations happen within the first 300 or so steps, after which the curves flatten (Fig.~\ref{fig:survival});
pure SCP's keeps falling for longer.
 in the appendix. Files and numbers as in experiments_results.tex.
% Order is by importance; drop subsections from the end if the appendix must be short.
% =====================================================================================================================
\providecommand{\figdir}{../paper_graphs}

\section{Additional Results}\label{app:results}

% ---------------------------------------------------------------------------------------------------------------------
\subsection{Per-model results}\label{app:per_model}

\begin{table}[htbp]
  \centering
  \caption{Table~\ref{tab:pooled} on the $c=1.2$ world models.}
  \label{tab:pooled_c1_2}
  \small
  \input{\figdir/appendix/table_pooled_c1_2}
\end{table}

\begin{table}[htbp]
  \centering
  \caption{Table~\ref{tab:pooled} on the unbounded ($c=\infty$) world models.}
  \label{tab:pooled_cinf}
  \small
  \input{\figdir/appendix/table_pooled_cinf}
\end{table}

Table~\ref{tab:per_model} gives each planner's violation rate and return on every $N=10$ world model, pooled over
$\epsilon\ge0.1$, and Figs.~\ref{fig:grid_violations} and~\ref{fig:grid_return} show the full sweeps. Pure SCP is
significantly less safe than SCP+MPPI on every model. Against the samplers, the only significant difference is MPPI
with 512 sequences on PointGoal2 with $c=1.2$, which violates less.

\begin{table}[htbp]
  \centering
  \caption{Each planner on each $N=10$ world model, pooled over $\epsilon\ge0.1$: violation rate (Wilson) and return
  (bootstrap) with 95\% intervals, and McNemar's $p$ for SCP+MPPI against it on paired seeds, Holm-corrected over the
  six models.}
  \label{tab:per_model}
  \small
  \input{\figdir/appendix/table_per_model}
\end{table}

\begin{figure}[htbp]
  \centering
  \includegraphics[width=\linewidth]{\figdir/appendix/fig_eps_sweep_all_bounds_violations.pdf}
  \caption{Violation rate against $\epsilon$ for every planner on every $N=10$ world model (rows: tasks; columns:
  Lipschitz bound $c$), with 95\% Wilson intervals.}
  \label{fig:grid_violations}
\end{figure}

\begin{figure}[htbp]
  \centering
  \includegraphics[width=\linewidth]{\figdir/appendix/fig_eps_sweep_all_bounds_return.pdf}
  \caption{Return against $\epsilon$ for every planner on every $N=10$ world model, with 95\% bootstrap intervals.}
  \label{fig:grid_return}
\end{figure}

% ---------------------------------------------------------------------------------------------------------------------
\subsection{Lipschitz bound and SCP's convergence}\label{app:lipschitz}

\begin{figure}[htbp]
  \centering
  \includegraphics[width=\linewidth]{\figdir/appendix/fig_lipschitz_convergence.pdf}
  \caption{Share of SCP+MPPI's steps that reached the 40-iteration cap against $\epsilon$, for each Lipschitz bound,
  at $N=10$.}
  \label{fig:lipschitz_convergence}
\end{figure}

Without a Lipschitz bound, SCP+MPPI reaches the iteration cap on more steps than with either bound at every margin,
except at $\epsilon=0.5$ on PointGoal2, where $c=1.5$ and $c=\infty$ are within 0.02 (Fig.~\ref{fig:lipschitz_convergence}).
Which bound converges best differs by task: $c=1.5$ on PointGoal1, $c=1.2$ on PointGoal2.

% ---------------------------------------------------------------------------------------------------------------------
\subsection{What more samples buy at \texorpdfstring{$N=10$}{N=10}}\label{app:samples}

\begin{figure}[htbp]
  \centering
  \includegraphics[width=\linewidth]{\figdir/ablation/fig_sample_vs_time_matched.pdf}
  \caption{Fig.~\ref{fig:eps_sweep} with the samplers at both budgets: 512 sequences (open markers, dashed) and
  time-matched (filled, solid).}
  \label{fig:sample_vs_time_matched}
\end{figure}

\begin{figure}[htbp]
  \centering
  \includegraphics[width=\linewidth]{\figdir/ablation/fig_return_vs_violations.pdf}
  \caption{Return against violation rate on each task's $c=1.5$ model at $N=10$, pooled over $\epsilon\ge0.1$, with
  95\% intervals. Arrows lead from each sampler with 512 sequences to the same sampler time-matched.}
  \label{fig:trade_off}
\end{figure}

Figs.~\ref{fig:sample_vs_time_matched} and~\ref{fig:trade_off} show what 128 times as many sequences buy the samplers.
On PointGoal1's $c=1.5$ model, the time-matched samplers earn more than with 512 sequences but violate more often
(MPPI: return 0.78 to 1.15, violation rate 0.19 to 0.24; CEM: 0.84 to 0.94 and 0.14 to 0.26), while SCP+MPPI reaches
the higher return at the lower violation rate (1.15 and 0.17). On PointGoal2 the extra sequences raise return (0.35 to
0.51 and 0.49) without changing the violation rate. These shifts lie within the intervals; the paired tests of
Sec.~\ref{sec:q1} find no significant difference in safety between SCP+MPPI and either budget.

% ---------------------------------------------------------------------------------------------------------------------
\subsection{All budgets at every horizon}\label{app:horizon_budgets}

\begin{figure}[htbp]
  \centering
  \includegraphics[width=\linewidth]{\figdir/ablation/fig_horizon_all_budgets.pdf}
  \caption{Fig.~\ref{fig:horizon} with the samplers at both 512 sequences (open markers, dashed) and their
  time-matched budget (filled, solid).}
  \label{fig:horizon_all}
\end{figure}

With 512 sequences the samplers' return at $N=20$ is close to SCP+MPPI's (MPPI 2.07, CEM 2.58, against 2.34; neither
difference significant), and at $N=30$ below it (1.59 and 1.22 against 2.26; $p=0.043$ and $0.013$)
(Fig.~\ref{fig:horizon_all}). Together with Fig.~\ref{fig:compute}, this shows that the samplers' best budget depends on
the sampler and the horizon.

% ---------------------------------------------------------------------------------------------------------------------
\subsection{Does optimisation exploit the model?}\label{app:optimism}

\begin{figure}[htbp]
  \centering
  \includegraphics[width=\linewidth]{\figdir/ablation/fig_model_optimism.pdf}
  \caption{Model optimism: each executed plan's predicted return (discounted predicted rewards plus terminal value)
  minus the discounted return received from that step on, averaged over each episode, with 95\% bootstrap intervals.
  $N=10$ pools every world model ($\epsilon\ge0.1$); $N=20$ and $30$ are the horizon study's models.}
  \label{fig:optimism}
\end{figure}

A gradient-based optimiser might exploit model error, finding plans whose predicted return the environment does not
deliver. We measure each executed plan's optimism: its predicted return, whose terminal value carries the goal bonus,
minus the discounted return actually received from that step on, using steps at least 299 steps before the episode
ends so that the part of the future the episode did not reach weighs under 5\%. Every planner is optimistic by about
2, as expected of a value fitted to the 0.7 expectile, but by nearly the same amount (Fig.~\ref{fig:optimism}): 2.06
(pure SCP) to 2.19 at $N=10$, 2.20 to 2.44 at $N=20$ and 2.41 to 2.56 at $N=30$, with SCP+MPPI inside the samplers'
range at every horizon and no trend with the sampling budget. We find no evidence that SCP, or a larger sampling
budget, exploits the model more than sampling with 512 sequences does.

% ---------------------------------------------------------------------------------------------------------------------
\subsection{Warned and unwarned violations}\label{app:crashes}

\begin{figure}[htbp]
  \centering
  \includegraphics[width=\linewidth]{\figdir/ablation/fig_crash_types.pdf}
  \caption{Violations per 100 episodes over the six $N=10$ models and $\epsilon\ge0.1$, split into those preceded
  within $N$ steps by a plan that broke the barrier in the model (warned, light) and those that were not (unwarned,
  dark, counts above the bars).}
  \label{fig:crash_types}
\end{figure}

Few violations come without warning (Fig.~\ref{fig:crash_types}). SCP+MPPI has 14 unwarned violations in 880
episodes, against 6 and 7 for the 512-sequence samplers and 7 and 13 for the time-matched ones. That stronger
optimisation, by SCP or by more samples, leads to more violations the model did not foresee would be consistent with
it planning closer to where the model is wrong, but paired by seed none of these differences is significant
($p\ge0.12$).

% ---------------------------------------------------------------------------------------------------------------------
\subsection{When violations happen}\label{app:survival}

\begin{figure}[htbp]
  \centering
  \includegraphics[width=\linewidth]{\figdir/appendix/fig_survival_pg1_c1_5.pdf}
  \caption{Share of episodes without a violation by each step, at each margin, on PointGoal1's $c=1.5$ model at
  $N=10$. Episodes end early only at a violation, so these are Kaplan--Meier curves without censoring.}
  \label{fig:survival}
\end{figure}

Most violations happen within the first 300 or so steps, after which the curves flatten (Fig.~\ref{fig:survival});
pure SCP's keeps falling for longer.
. Citations: experiments_refs.bib.
%
% Every figure and table is written by the repository's scripts, run from the repository root:
%   python src/scripts/analyse_results.py  -> src/save_data/experiments/analysis_results/*.csv
%   python src/scripts/paper_figures.py    -> paper_graphs/{main_text,ablation,appendix}/ (figures and both tables)
%   python src/scripts/env_figure.py       -> paper_graphs/main_text/fig_environments.pdf
% \figdir is the paper_graphs folder; define it before \input if the paper lives elsewhere (e.g. on Overleaf).
%
% Notation follows the method: N horizon, \epsilon barrier margin, \alpha DCBF rate, c Lipschitz bound, \tau risk
% threshold, g_\Phi margin head. Labels of the method's parts used here (rename to match yours): sec:regularity,
% sec:cbf, sec:scp, sec:init, eq:ocp.
% Every number is from the analysis of the runs on main (2d4240b). "% TODO" marks a fact to confirm before submission.
% =====================================================================================================================
\providecommand{\figdir}{../paper_graphs}

\section{Experimental Setup}\label{sec:setup}

Our experiments answer four questions.
\textbf{(Q1)}~Does SCP, started from a sampled plan, plan better than sampling-based MPC given the same wall-clock time?
\textbf{(Q2)}~How does the comparison change as the horizon, and with it the number of decision variables, grows?
\textbf{(Q3)}~What does SCP's convergence tell us online?
\textbf{(Q4)}~Does meeting the barrier in the learned model keep the robot safe in the environment, and what do the
regularity constraints cost?

% ---------------------------------------------------------------------------------------------------------------------
\subsection{Environments}\label{sec:envs}

\begin{figure}[t]
  \centering
  \includegraphics[width=0.62\linewidth]{\figdir/main_text/fig_environments.pdf}
  \caption{The two tasks from above, one layout each (layouts are drawn from the episode's seed). The point robot
  (red) must reach the goal (green) without entering a hazard (blue); in PointGoal2, touching a vase (cyan) is also a
  violation.}
  \label{fig:envs}
\end{figure}

We use \texttt{SafetyPointGoal1} and \texttt{SafetyPointGoal2} from Safety-Gymnasium~\citep{ji2023safety}
(Fig.~\ref{fig:envs}). A point robot with two actuators (forward force and turning, $u\in[-1,1]^2$) must reach a goal
of radius $0.3$\,m, which moves to a new position each time it is reached, within $T=1000$ steps. The reward is the
per-step reduction in distance to the goal plus a bonus of 1 for reaching it. The robot observes 60 quantities: 12
proprioceptive signals (accelerometer, velocimeter, gyroscope and magnetometer) and three 16-ray pseudo-lidars measuring
the proximity of the goal, the hazards and the vases. Level~1 places 8 hazards of radius $0.2$\,m and one cost-free
vase in a $3\times3$\,m region; level~2 places 10 hazards and 10 vases in a $4\times4$\,m region and also charges for
touching or moving a vase. A step is a \emph{violation} if it incurs any cost. We treat violations as absorbing
failures: an episode ends at its first violation, so the \emph{violation rate} is the share of episodes that fail and
the \emph{return} is the reward collected before failure or the time limit.

These tasks suit our question for four reasons.
(i)~DSRL~\citep{liu2023datasets} provides offline datasets for both, so the world model is learned without any
interaction, which is the setting of this paper.
(ii)~The constraint is geometric (stay out of regions of the plane) and observed through the lidar, so the barrier's
margin is a distance in metres, the margin $\epsilon$ has a physical meaning, and safety in the model can be checked
against safety in the environment (Sec.~\ref{sec:latent_safety}).
(iii)~Level~2 is more cluttered (20 obstacles in 16\,m$^2$ against 8 in 9\,m$^2$) and adds a second, movable kind of
obstacle, so the two levels, together with $\epsilon$, vary how restrictive the safe set is.
(iv)~The action is two-dimensional, which favours sampling: at $N=10$ a plan has only 20 decision variables, so the
comparison is a demanding one for SCP. The horizon study (Sec.~\ref{sec:horizon}) raises this to 40 and 60.

% ---------------------------------------------------------------------------------------------------------------------
\subsection{Offline data and world models}\label{sec:wm_setup}

For each task we train on the corresponding DSRL dataset (2,022 trajectories for PointGoal1 and 3,442 for PointGoal2,
of up to 1000 steps each), truncating every trajectory at its first violation to match the absorbing failure above.
This leaves about 0.87\,M transitions for PointGoal1 and 1.15\,M for PointGoal2; only 329 of PointGoal2's trajectories
run their full length without a violation.
% TODO: PG1's exact transition count (FilterDataset.ipynb prints it for PG2 only: 1,145,452); one sentence on how DSRL
% collected the data (behaviour policies, cost budgets).
The encoder drops seven proprioceptive channels (the accelerometer and the out-of-plane components of velocity,
angular rate and magnetic field) and maps the remaining 53 to a latent state of $d_z=32$ dimensions. The four networks
the planner differentiates through (dynamics, reward, value and margin) are MLPs with two hidden layers of 256 units,
each with Lipschitz constant at most $c$ (Sec.~\ref{sec:regularity}). We train one model per task and bound,
$c\in\{1.2,1.5,\infty\}$ ($c=\infty$: unbounded), for $N=10$, and two more PointGoal1 models with $c=1.5$ for $N=20$
and $30$. Each is trained for 250 epochs of 1000 (PointGoal1) or 1500 (PointGoal2) minibatches of 256 windows of $N+1$
consecutive transitions, with learning rate $3\times10^{-4}$, the loss at rollout step $t$ weighted by $\rho^t$
($\rho=0.9$, $0.95$, $0.99$ for $N=10$, $20$, $30$), discount $\gamma=0.99$ and value expectile $0.7$. The margin head
$g_\Phi$ regresses the distance from the robot to the nearest hazard edge (in PointGoal2, hazard or vase edge),
computed from the lidar, clipped at 1\,m, with samples near an edge up-weighted.
% TODO: optimiser (Adam?); state that each model is a single training seed (limitations).

% ---------------------------------------------------------------------------------------------------------------------
\subsection{Planners}\label{sec:planners}

All planners solve the same receding-horizon problem~\eqref{eq:ocp} in the same model: maximise
$\sum_{k=0}^{N-1}\gamma^k r(z_k,u_k)+\gamma^N V(z_N)-\rho_u\|u\|^2$ ($\rho_u=10^{-4}$) subject to
$(1-\alpha)h(z_k)-h(z_{k+1})\le\tau=0$ at every step, with $h(z)=g_\Phi(z)-\epsilon$ and $\alpha=0.05$. Each executes
the first action of its plan and replans at every step.

\textbf{SCP} is the method of Sec.~\ref{sec:scp} started from its previous plan shifted by one step (zeros at the first
step): at most 40 prox-linear iterations, each a convex QP solved by Clarabel~\citep{goulart2024clarabel}, with exact
$\ell_1$ penalties $w_{\mathrm{dyn}}=10^3$ and $w_{\mathrm{safe}}=10^2$, a trust-region weight that starts at 1, is
floored at $10^{-2}$ and is carried between steps, acceptance thresholds $(0.25,0.75)$, and a stop when the
prox-gradient norm falls below $10^{-2}$.
\textbf{SCP+MPPI}, our method, first runs MPPI with 512 sequences (below) and starts SCP from whichever of MPPI's plan
and its own shifted plan has the lower penalty merit (Sec.~\ref{sec:init}). SCP accepts only steps that lower the
merit, so the plan SCP+MPPI executes is never worse under the merit than MPPI's, whether or not SCP converges.
\textbf{MPPI}~\citep{williams2017information} and \textbf{CEM}~\citep{rubinstein1999cross} run 6 iterations of $J$
action sequences drawn from a Gaussian that is refitted to the best $J/8$ (weighted by $\exp(\text{score}/0.5)$ in MPPI,
uniformly in CEM) and warm-started from the previous mean shifted by one step. A sequence whose rollout breaks the
barrier at any step ranks below every sequence that meets it; if none meets it, sequences are ranked by their total
barrier excess, the quantity SCP's $\ell_1$ penalty charges. The executed plan is the refitted mean if its rollout
meets the barrier, and otherwise the best sequence sampled.

We run the samplers with $J=512$, the proposal SCP+MPPI uses, and \emph{time-matched}, with $J$ chosen so that a step
takes about as long as SCP+MPPI's: $J=65{,}536$ at $N=10$ (about 0.70\,s against 0.40--0.71\,s for SCP+MPPI, timed on
the same idle GPU), and $95{,}040$ and $120{,}000$ at $N=20$ and $30$ (mean logged plan times 2.04\,s and 3.81\,s,
against 2.16\,s and 3.61\,s for SCP+MPPI). Each run used one 1g.10gb slice of an NVIDIA A100, with the QP solved on the
CPU.
% TODO: the 65,536 time is interpolated between 32,768 (0.33 s) and 131,072 (1.43 s) sequences; time it directly or
% say so. Confirm the hardware sentence.

% ---------------------------------------------------------------------------------------------------------------------
\subsection{Protocol and statistics}\label{sec:protocol}

We sweep the margin $\epsilon\in\{0,0.1,0.2,0.3,0.5\}$\,m on every $N=10$ model and fix $\epsilon=0.2$\,m for the
horizon study. Each configuration runs on seeds $0$--$39$, or $0$--$29$ for the time-matched samplers, the horizon
study, PointGoal2's $c=\infty$ model and SCP+MPPI on PointGoal1's $c=\infty$ model. The seed fixes the layout, the
initial state and the planner's random numbers, so every comparison between planners or models is paired by seed. We
report violation rates with Wilson 95\% intervals and mean returns with bootstrap 95\% intervals ($10^4$ resamples).
Paired comparisons report the difference in violation rate and in return with a 95\% bootstrap interval, and the
probability, given the data, that SCP+MPPI's violation rate is lower or its return higher than the baseline's, from a
Bayesian bootstrap~\citep{rubin1981bayesian}. When pooling over margins and tasks we use $\epsilon\ge0.1$, since at
$\epsilon=0$ every planner fails often (Sec.~\ref{sec:q1}), and resample whole seeds, since a seed's episodes share a
layout. The per-model results in the appendix use McNemar's exact test on violations, Holm-corrected over the six models,
and the Wilcoxon signed-rank test on returns.

% =====================================================================================================================
\section{Results}\label{sec:results}

% ---------------------------------------------------------------------------------------------------------------------
\subsection{SCP needs a sampled start, and then matches time-matched sampling at \texorpdfstring{$N=10$}{N=10}}\label{sec:q1}

\begin{figure}[t]
  \centering
  \includegraphics[width=\linewidth]{\figdir/main_text/fig_eps_sweep_c1_5.pdf}
  \caption{Violation rate (top) and return (bottom) against the margin $\epsilon$ on each task's $c=1.5$ model at
  $N=10$, with 95\% intervals; 40 episodes per point (30 for the time-matched samplers).}
  \label{fig:eps_sweep}
\end{figure}

\begin{table}[t]
  \centering
  \caption{SCP+MPPI minus each baseline on the $c=1.5$ world models at $N=10$, pooled over both tasks and
  $\epsilon\ge0.1$ and paired by task, margin and seed (Tables~\ref{tab:pooled_c1_2} and~\ref{tab:pooled_cinf} in the
  appendix give the other bounds). A negative $\Delta$ violation rate or a positive $\Delta$ return favours SCP+MPPI.
  Brackets: 95\% intervals resampling whole seeds. $P$(safer) and $P$(higher return): the probability, given the data,
  that SCP+MPPI's violation rate is lower, or its return higher, than the baseline's on these models (Bayesian bootstrap
  over seeds).}
  \label{tab:pooled}
  \small
  \input{\figdir/main_text/table_pooled_c1_5}
\end{table}

\textbf{SCP needs a sampled start.} Started from its own shifted plan, SCP fails in 55\% of episodes over all $N=10$
models and $\epsilon\ge0.1$, against 25\% for SCP+MPPI on the same layouts (Table~\ref{tab:pooled}: $-29.8$ points,
$p=2\times10^{-43}$), at the same return. Pure SCP does not fail for want of convergence: on PointGoal1's $c=1.5$ model
it certifies nearly as many steps as SCP+MPPI (62--95\% against 63--97\% across $\epsilon$; Sec.~\ref{sec:certificates}). It fails
as a local method does, converging to the local solution nearest a poor start: it converges with barrier slack in use,
meaning no safe plan nearby, on more steps (2--19\% against 0--11\% on PointGoal1) and finds a plan that meets the
barrier in the model less often (88\% of steps against 92\%). On the $c=1.5$ models, SCP+MPPI starts from MPPI's plan on
33--43\% of steps on PointGoal1 and 38--55\% on PointGoal2 ($\epsilon\ge0.1$).

\textbf{With a sampled start, SCP+MPPI is as safe as sampling and earns more than its proposal.} Against MPPI and CEM
with the 512 sequences that seed it, SCP+MPPI's violation rate differs by $+1.9$ and $+2.6$ points ($p=0.24$ and
$0.098$), while its return is higher by 0.21 and 0.18 with intervals that exclude zero. Against samplers given the same
time, with 128 times as many sequences, the differences in violations ($-0.1$ and $+0.7$ points) and return ($+0.06$
and $+0.09$) are small and not significant. At $N=10$, then, SCP's refinement adds about as much return as 128 times
more samples, and neither makes the planner safer: with a two-dimensional action, sampling finds plans that meet the
barrier about as often as optimisation does (Sec.~\ref{sec:latent_safety}). The per-model results agree
(Table~\ref{tab:per_model}): no time-matched sampler differs significantly from SCP+MPPI in violations on any model,
and the one significant difference against the 512-sequence samplers favours MPPI on PointGoal2 with $c=1.2$ (23\%
against 34\%, Holm-corrected $p=0.031$).

\textbf{The margin is the main safety lever.} Without a margin ($\epsilon=0$) every planner fails in 30--63\% of
episodes, and pure SCP in 72--85\% (Fig.~\ref{fig:eps_sweep}): the barrier holds only in the model, so a margin must
absorb the model's error. A margin of 0.1--0.3\,m cuts SCP+MPPI's violation rate from 40--50\% to 15--23\% at a moderate
cost in return. At $\epsilon=0.5$\,m the planners turn conservative and return collapses (SCP+MPPI on PointGoal1: 0.38
against 1.75 at $\epsilon=0.1$), and on the more cluttered PointGoal2 violations rise again for every planner
(40--45\%; pure SCP 70\%) as the barrier becomes infeasible in many states, which SCP reports (Sec.~\ref{sec:certificates}).

% ---------------------------------------------------------------------------------------------------------------------
\subsection{At longer horizons SCP+MPPI earns more than time-matched sampling}\label{sec:horizon}

\begin{figure}[t]
  \centering
  \includegraphics[width=\linewidth]{\figdir/main_text/fig_horizon.pdf}
  \caption{Violation rate and return against the horizon $N$ on PointGoal1 ($c=1.5$, $\epsilon=0.2$), each sampler at
  its horizon's time-matched budget, with 95\% intervals. Each $N$ uses a world model trained on windows of $N+1$
  steps; 30 episodes per point (40 for SCP+MPPI and SCP at $N=10$).}
  \label{fig:horizon}
\end{figure}

% CUT 1 if over the page limit: move this figure to the appendix and keep the paragraph's numbers.
\begin{figure}[t]
  \centering
  \includegraphics[width=\linewidth]{\figdir/main_text/fig_return_vs_compute.pdf}
  \caption{Return against mean plan time per step at $N=20$ and $30$, with 95\% intervals. Each sampler is shown at
  every budget we ran (sequences per iteration along the bottom); SCP+MPPI and SCP are single points.}
  \label{fig:compute}
\end{figure}

We repeat the comparison on PointGoal1 ($c=1.5$, $\epsilon=0.2$) at $N=20$ and $30$, re-matching the samplers' budgets
in time at each horizon (Fig.~\ref{fig:horizon}). At both horizons SCP+MPPI earns more than the time-matched samplers,
at the same safety (2 and 1 violations in 30 episodes, against 3 and 0). Its return is higher by 0.65 [0.03, 1.30]
(MPPI, $p=0.021$) and 0.93 [0.15, 1.69] (CEM, $p=0.009$) at $N=20$, and by 1.13 [0.41, 1.93] ($p<0.001$) and 0.98
[0.08, 1.93] ($p=0.005$) at $N=30$.

The samplers do not turn extra time into return (Fig.~\ref{fig:compute}). Their largest budgets earn less than smaller
ones at both horizons, and their best budget changes with the sampler and the horizon (512 sequences for CEM at
$N=20$, 32,768 at $N=30$). At its best budget each sampler ties with SCP+MPPI (CEM: $-0.24$ at $N=20$, $+0.11$ at
$N=30$; neither significant), but that budget is not known in advance. SCP+MPPI used the same settings at every
horizon, and its refinement adds 0.27 (not significant) and 0.67 ($p=0.043$) to the return of the 512-sequence MPPI
plan it starts from.
% Possible reason for the samplers' trend, not tested: the elite set grows with J (top-k = J/8), so a larger budget
% refits to the same top eighth of the samples instead of concentrating on the best of them.

Two caveats apply. The horizon study uses one world model per horizon, one margin and 30 seeds, and the model changes
with $N$, so returns compare within a horizon but not across horizons. And within its 40-iteration cap SCP+MPPI
converged on 53\% of steps at $N=20$ and 32\% at $N=30$, executing its last accepted iterate on the rest: the
comparison is at matched time, not at convergence.

% ---------------------------------------------------------------------------------------------------------------------
\subsection{SCP's convergence certifies the plan and detects infeasibility}\label{sec:certificates}

\begin{figure}[t]
  \centering
  \includegraphics[width=\linewidth]{\figdir/main_text/fig_scp_stop_reasons.pdf}
  \caption{How SCP+MPPI's steps ended on each task's $c=1.5$ model at $N=10$: converged with the plan's rollout meeting
  the barrier (certified), converged with barrier slack in use, or at the 40-iteration cap.}
  \label{fig:stops}
\end{figure}

Each SCP step ends in one of three ways (Fig.~\ref{fig:stops}). Because SCP's reference trajectory is always the model
rollout of its plan, a step that stops as converged, with prox-gradient norm below $10^{-2}$ and barrier excess below
$10^{-4}$, executes an approximately stationary plan that meets the barrier at every step in the model: convergence
itself certifies the plan, at no extra cost. This is the most common outcome. On the $c=1.5$ models SCP+MPPI certifies
97\% of steps at $\epsilon=0$, falling to 63\% at $\epsilon=0.5$, on PointGoal1, and 83\% falling to 60\% on
PointGoal2, using 9--22 iterations on average.

A step that converges with barrier slack in use signals something different: a stationary point at which the barrier
cannot be met, that is, no safe plan near the current one. Its share grows with the margin, from 0 to 11\% of steps on
PointGoal1 and from 1\% to 22\% on PointGoal2, where at $\epsilon=0.5$ it coincides with the rise in violations
(Sec.~\ref{sec:q1}). Sampling planners give no such signal: when no sampled sequence meets the barrier they cannot tell
an empty safe set from too few samples. The remaining steps reach the iteration cap, a share that grows with $\epsilon$
(3\% to 27\% on PointGoal1) and with the horizon (46\% at $N=20$ and 68\% at $N=30$ on the same task and model bound).

% ---------------------------------------------------------------------------------------------------------------------
\subsection{Safety in the model predicts safety in the environment}\label{sec:latent_safety}

% CUT 2 if over the page limit: move this figure to the appendix; the paragraph stands on its numbers.
\begin{figure}[t]
  \centering
  \includegraphics[width=\linewidth]{\figdir/main_text/fig_latent_safety.pdf}
  \caption{Share of steps followed by a violation within $N$ steps, after a plan whose model rollout met the barrier
  and after one whose rollout did not, over the six $N=10$ models and $\epsilon\ge0.1$, with 95\% Wilson intervals
  (log scale).}
  \label{fig:latent_safety}
\end{figure}

The barrier is enforced only in the learned model; Fig.~\ref{fig:latent_safety} tests whether that is enough. After a
plan that meets the barrier in the model, a violation follows within $N$ steps on $2$--$10\times10^{-4}$ of steps,
whichever planner made the plan (including the 512-sequence samplers); after a plan that does not, on
$3$--$7\times10^{-2}$, 56 to 131 times as often. Conversely, 93--98\% of violations were preceded within $N$ steps by a
plan that broke the barrier in the model: the model nearly always warns before a violation. Violations therefore arise
mainly in states where the planner finds no plan that is safe in the model, not from plans that are safe in the model
but not in the environment. All planners find such a plan on a similar share of steps (90--92\%; pure SCP 88\%), which
is why optimising harder does not make the planner safer at $N=10$ (Sec.~\ref{sec:q1}).

% ---------------------------------------------------------------------------------------------------------------------
\subsection{Regularity constraints}\label{sec:regularity_results}

Bounding the networks' Lipschitz constants at $c=1.5$ costs no measurable safety or return: against the unbounded
models, no planner's violation rate or return differs significantly on either task (McNemar $p\ge0.18$; return
intervals include zero;
Figs.~\ref{fig:grid_violations} and~\ref{fig:grid_return} in the appendix). It does make SCP converge more reliably
(Fig.~\ref{fig:lipschitz_convergence}). At $\epsilon$ from 0.1 to 0.3, SCP+MPPI reaches the iteration cap on 6--16\% of
steps on PointGoal1 and 16--21\% on PointGoal2 with $c=1.5$, against 33--44\% and 31--35\% unbounded, needs fewer
iterations (12--16 against 24--28 on PointGoal1), and so certifies more of its plans (78--92\% against 55--65\% on
PointGoal1; 74--83\% against 53--62\% on PointGoal2). The tighter bound $c=1.2$ makes the model less conservative:
returns are higher, but so are violations, significantly on PointGoal1 for SCP+MPPI (44 against 28 of 160 episodes,
$p=0.044$), MPPI ($p=0.009$) and CEM ($p=0.005$), and on PointGoal2 for SCP+MPPI (54 against 38, $p=0.036$). We
therefore use $c=1.5$ for the main results.
. Appendix: % =====================================================================================================================
% Appendix results: % =====================================================================================================================
% Appendix results: % =====================================================================================================================
% Appendix results: \input{appendix_results} in the appendix. Files and numbers as in experiments_results.tex.
% Order is by importance; drop subsections from the end if the appendix must be short.
% =====================================================================================================================
\providecommand{\figdir}{../paper_graphs}

\section{Additional Results}\label{app:results}

% ---------------------------------------------------------------------------------------------------------------------
\subsection{Per-model results}\label{app:per_model}

\begin{table}[htbp]
  \centering
  \caption{Table~\ref{tab:pooled} on the $c=1.2$ world models.}
  \label{tab:pooled_c1_2}
  \small
  \input{\figdir/appendix/table_pooled_c1_2}
\end{table}

\begin{table}[htbp]
  \centering
  \caption{Table~\ref{tab:pooled} on the unbounded ($c=\infty$) world models.}
  \label{tab:pooled_cinf}
  \small
  \input{\figdir/appendix/table_pooled_cinf}
\end{table}

Table~\ref{tab:per_model} gives each planner's violation rate and return on every $N=10$ world model, pooled over
$\epsilon\ge0.1$, and Figs.~\ref{fig:grid_violations} and~\ref{fig:grid_return} show the full sweeps. Pure SCP is
significantly less safe than SCP+MPPI on every model. Against the samplers, the only significant difference is MPPI
with 512 sequences on PointGoal2 with $c=1.2$, which violates less.

\begin{table}[htbp]
  \centering
  \caption{Each planner on each $N=10$ world model, pooled over $\epsilon\ge0.1$: violation rate (Wilson) and return
  (bootstrap) with 95\% intervals, and McNemar's $p$ for SCP+MPPI against it on paired seeds, Holm-corrected over the
  six models.}
  \label{tab:per_model}
  \small
  \input{\figdir/appendix/table_per_model}
\end{table}

\begin{figure}[htbp]
  \centering
  \includegraphics[width=\linewidth]{\figdir/appendix/fig_eps_sweep_all_bounds_violations.pdf}
  \caption{Violation rate against $\epsilon$ for every planner on every $N=10$ world model (rows: tasks; columns:
  Lipschitz bound $c$), with 95\% Wilson intervals.}
  \label{fig:grid_violations}
\end{figure}

\begin{figure}[htbp]
  \centering
  \includegraphics[width=\linewidth]{\figdir/appendix/fig_eps_sweep_all_bounds_return.pdf}
  \caption{Return against $\epsilon$ for every planner on every $N=10$ world model, with 95\% bootstrap intervals.}
  \label{fig:grid_return}
\end{figure}

% ---------------------------------------------------------------------------------------------------------------------
\subsection{Lipschitz bound and SCP's convergence}\label{app:lipschitz}

\begin{figure}[htbp]
  \centering
  \includegraphics[width=\linewidth]{\figdir/appendix/fig_lipschitz_convergence.pdf}
  \caption{Share of SCP+MPPI's steps that reached the 40-iteration cap against $\epsilon$, for each Lipschitz bound,
  at $N=10$.}
  \label{fig:lipschitz_convergence}
\end{figure}

Without a Lipschitz bound, SCP+MPPI reaches the iteration cap on more steps than with either bound at every margin,
except at $\epsilon=0.5$ on PointGoal2, where $c=1.5$ and $c=\infty$ are within 0.02 (Fig.~\ref{fig:lipschitz_convergence}).
Which bound converges best differs by task: $c=1.5$ on PointGoal1, $c=1.2$ on PointGoal2.

% ---------------------------------------------------------------------------------------------------------------------
\subsection{What more samples buy at \texorpdfstring{$N=10$}{N=10}}\label{app:samples}

\begin{figure}[htbp]
  \centering
  \includegraphics[width=\linewidth]{\figdir/ablation/fig_sample_vs_time_matched.pdf}
  \caption{Fig.~\ref{fig:eps_sweep} with the samplers at both budgets: 512 sequences (open markers, dashed) and
  time-matched (filled, solid).}
  \label{fig:sample_vs_time_matched}
\end{figure}

\begin{figure}[htbp]
  \centering
  \includegraphics[width=\linewidth]{\figdir/ablation/fig_return_vs_violations.pdf}
  \caption{Return against violation rate on each task's $c=1.5$ model at $N=10$, pooled over $\epsilon\ge0.1$, with
  95\% intervals. Arrows lead from each sampler with 512 sequences to the same sampler time-matched.}
  \label{fig:trade_off}
\end{figure}

Figs.~\ref{fig:sample_vs_time_matched} and~\ref{fig:trade_off} show what 128 times as many sequences buy the samplers.
On PointGoal1's $c=1.5$ model, the time-matched samplers earn more than with 512 sequences but violate more often
(MPPI: return 0.78 to 1.15, violation rate 0.19 to 0.24; CEM: 0.84 to 0.94 and 0.14 to 0.26), while SCP+MPPI reaches
the higher return at the lower violation rate (1.15 and 0.17). On PointGoal2 the extra sequences raise return (0.35 to
0.51 and 0.49) without changing the violation rate. These shifts lie within the intervals; the paired tests of
Sec.~\ref{sec:q1} find no significant difference in safety between SCP+MPPI and either budget.

% ---------------------------------------------------------------------------------------------------------------------
\subsection{All budgets at every horizon}\label{app:horizon_budgets}

\begin{figure}[htbp]
  \centering
  \includegraphics[width=\linewidth]{\figdir/ablation/fig_horizon_all_budgets.pdf}
  \caption{Fig.~\ref{fig:horizon} with the samplers at both 512 sequences (open markers, dashed) and their
  time-matched budget (filled, solid).}
  \label{fig:horizon_all}
\end{figure}

With 512 sequences the samplers' return at $N=20$ is close to SCP+MPPI's (MPPI 2.07, CEM 2.58, against 2.34; neither
difference significant), and at $N=30$ below it (1.59 and 1.22 against 2.26; $p=0.043$ and $0.013$)
(Fig.~\ref{fig:horizon_all}). Together with Fig.~\ref{fig:compute}, this shows that the samplers' best budget depends on
the sampler and the horizon.

% ---------------------------------------------------------------------------------------------------------------------
\subsection{Does optimisation exploit the model?}\label{app:optimism}

\begin{figure}[htbp]
  \centering
  \includegraphics[width=\linewidth]{\figdir/ablation/fig_model_optimism.pdf}
  \caption{Model optimism: each executed plan's predicted return (discounted predicted rewards plus terminal value)
  minus the discounted return received from that step on, averaged over each episode, with 95\% bootstrap intervals.
  $N=10$ pools every world model ($\epsilon\ge0.1$); $N=20$ and $30$ are the horizon study's models.}
  \label{fig:optimism}
\end{figure}

A gradient-based optimiser might exploit model error, finding plans whose predicted return the environment does not
deliver. We measure each executed plan's optimism: its predicted return, whose terminal value carries the goal bonus,
minus the discounted return actually received from that step on, using steps at least 299 steps before the episode
ends so that the part of the future the episode did not reach weighs under 5\%. Every planner is optimistic by about
2, as expected of a value fitted to the 0.7 expectile, but by nearly the same amount (Fig.~\ref{fig:optimism}): 2.06
(pure SCP) to 2.19 at $N=10$, 2.20 to 2.44 at $N=20$ and 2.41 to 2.56 at $N=30$, with SCP+MPPI inside the samplers'
range at every horizon and no trend with the sampling budget. We find no evidence that SCP, or a larger sampling
budget, exploits the model more than sampling with 512 sequences does.

% ---------------------------------------------------------------------------------------------------------------------
\subsection{Warned and unwarned violations}\label{app:crashes}

\begin{figure}[htbp]
  \centering
  \includegraphics[width=\linewidth]{\figdir/ablation/fig_crash_types.pdf}
  \caption{Violations per 100 episodes over the six $N=10$ models and $\epsilon\ge0.1$, split into those preceded
  within $N$ steps by a plan that broke the barrier in the model (warned, light) and those that were not (unwarned,
  dark, counts above the bars).}
  \label{fig:crash_types}
\end{figure}

Few violations come without warning (Fig.~\ref{fig:crash_types}). SCP+MPPI has 14 unwarned violations in 880
episodes, against 6 and 7 for the 512-sequence samplers and 7 and 13 for the time-matched ones. That stronger
optimisation, by SCP or by more samples, leads to more violations the model did not foresee would be consistent with
it planning closer to where the model is wrong, but paired by seed none of these differences is significant
($p\ge0.12$).

% ---------------------------------------------------------------------------------------------------------------------
\subsection{When violations happen}\label{app:survival}

\begin{figure}[htbp]
  \centering
  \includegraphics[width=\linewidth]{\figdir/appendix/fig_survival_pg1_c1_5.pdf}
  \caption{Share of episodes without a violation by each step, at each margin, on PointGoal1's $c=1.5$ model at
  $N=10$. Episodes end early only at a violation, so these are Kaplan--Meier curves without censoring.}
  \label{fig:survival}
\end{figure}

Most violations happen within the first 300 or so steps, after which the curves flatten (Fig.~\ref{fig:survival});
pure SCP's keeps falling for longer.
 in the appendix. Files and numbers as in experiments_results.tex.
% Order is by importance; drop subsections from the end if the appendix must be short.
% =====================================================================================================================
\providecommand{\figdir}{../paper_graphs}

\section{Additional Results}\label{app:results}

% ---------------------------------------------------------------------------------------------------------------------
\subsection{Per-model results}\label{app:per_model}

\begin{table}[htbp]
  \centering
  \caption{Table~\ref{tab:pooled} on the $c=1.2$ world models.}
  \label{tab:pooled_c1_2}
  \small
  \input{\figdir/appendix/table_pooled_c1_2}
\end{table}

\begin{table}[htbp]
  \centering
  \caption{Table~\ref{tab:pooled} on the unbounded ($c=\infty$) world models.}
  \label{tab:pooled_cinf}
  \small
  \input{\figdir/appendix/table_pooled_cinf}
\end{table}

Table~\ref{tab:per_model} gives each planner's violation rate and return on every $N=10$ world model, pooled over
$\epsilon\ge0.1$, and Figs.~\ref{fig:grid_violations} and~\ref{fig:grid_return} show the full sweeps. Pure SCP is
significantly less safe than SCP+MPPI on every model. Against the samplers, the only significant difference is MPPI
with 512 sequences on PointGoal2 with $c=1.2$, which violates less.

\begin{table}[htbp]
  \centering
  \caption{Each planner on each $N=10$ world model, pooled over $\epsilon\ge0.1$: violation rate (Wilson) and return
  (bootstrap) with 95\% intervals, and McNemar's $p$ for SCP+MPPI against it on paired seeds, Holm-corrected over the
  six models.}
  \label{tab:per_model}
  \small
  \input{\figdir/appendix/table_per_model}
\end{table}

\begin{figure}[htbp]
  \centering
  \includegraphics[width=\linewidth]{\figdir/appendix/fig_eps_sweep_all_bounds_violations.pdf}
  \caption{Violation rate against $\epsilon$ for every planner on every $N=10$ world model (rows: tasks; columns:
  Lipschitz bound $c$), with 95\% Wilson intervals.}
  \label{fig:grid_violations}
\end{figure}

\begin{figure}[htbp]
  \centering
  \includegraphics[width=\linewidth]{\figdir/appendix/fig_eps_sweep_all_bounds_return.pdf}
  \caption{Return against $\epsilon$ for every planner on every $N=10$ world model, with 95\% bootstrap intervals.}
  \label{fig:grid_return}
\end{figure}

% ---------------------------------------------------------------------------------------------------------------------
\subsection{Lipschitz bound and SCP's convergence}\label{app:lipschitz}

\begin{figure}[htbp]
  \centering
  \includegraphics[width=\linewidth]{\figdir/appendix/fig_lipschitz_convergence.pdf}
  \caption{Share of SCP+MPPI's steps that reached the 40-iteration cap against $\epsilon$, for each Lipschitz bound,
  at $N=10$.}
  \label{fig:lipschitz_convergence}
\end{figure}

Without a Lipschitz bound, SCP+MPPI reaches the iteration cap on more steps than with either bound at every margin,
except at $\epsilon=0.5$ on PointGoal2, where $c=1.5$ and $c=\infty$ are within 0.02 (Fig.~\ref{fig:lipschitz_convergence}).
Which bound converges best differs by task: $c=1.5$ on PointGoal1, $c=1.2$ on PointGoal2.

% ---------------------------------------------------------------------------------------------------------------------
\subsection{What more samples buy at \texorpdfstring{$N=10$}{N=10}}\label{app:samples}

\begin{figure}[htbp]
  \centering
  \includegraphics[width=\linewidth]{\figdir/ablation/fig_sample_vs_time_matched.pdf}
  \caption{Fig.~\ref{fig:eps_sweep} with the samplers at both budgets: 512 sequences (open markers, dashed) and
  time-matched (filled, solid).}
  \label{fig:sample_vs_time_matched}
\end{figure}

\begin{figure}[htbp]
  \centering
  \includegraphics[width=\linewidth]{\figdir/ablation/fig_return_vs_violations.pdf}
  \caption{Return against violation rate on each task's $c=1.5$ model at $N=10$, pooled over $\epsilon\ge0.1$, with
  95\% intervals. Arrows lead from each sampler with 512 sequences to the same sampler time-matched.}
  \label{fig:trade_off}
\end{figure}

Figs.~\ref{fig:sample_vs_time_matched} and~\ref{fig:trade_off} show what 128 times as many sequences buy the samplers.
On PointGoal1's $c=1.5$ model, the time-matched samplers earn more than with 512 sequences but violate more often
(MPPI: return 0.78 to 1.15, violation rate 0.19 to 0.24; CEM: 0.84 to 0.94 and 0.14 to 0.26), while SCP+MPPI reaches
the higher return at the lower violation rate (1.15 and 0.17). On PointGoal2 the extra sequences raise return (0.35 to
0.51 and 0.49) without changing the violation rate. These shifts lie within the intervals; the paired tests of
Sec.~\ref{sec:q1} find no significant difference in safety between SCP+MPPI and either budget.

% ---------------------------------------------------------------------------------------------------------------------
\subsection{All budgets at every horizon}\label{app:horizon_budgets}

\begin{figure}[htbp]
  \centering
  \includegraphics[width=\linewidth]{\figdir/ablation/fig_horizon_all_budgets.pdf}
  \caption{Fig.~\ref{fig:horizon} with the samplers at both 512 sequences (open markers, dashed) and their
  time-matched budget (filled, solid).}
  \label{fig:horizon_all}
\end{figure}

With 512 sequences the samplers' return at $N=20$ is close to SCP+MPPI's (MPPI 2.07, CEM 2.58, against 2.34; neither
difference significant), and at $N=30$ below it (1.59 and 1.22 against 2.26; $p=0.043$ and $0.013$)
(Fig.~\ref{fig:horizon_all}). Together with Fig.~\ref{fig:compute}, this shows that the samplers' best budget depends on
the sampler and the horizon.

% ---------------------------------------------------------------------------------------------------------------------
\subsection{Does optimisation exploit the model?}\label{app:optimism}

\begin{figure}[htbp]
  \centering
  \includegraphics[width=\linewidth]{\figdir/ablation/fig_model_optimism.pdf}
  \caption{Model optimism: each executed plan's predicted return (discounted predicted rewards plus terminal value)
  minus the discounted return received from that step on, averaged over each episode, with 95\% bootstrap intervals.
  $N=10$ pools every world model ($\epsilon\ge0.1$); $N=20$ and $30$ are the horizon study's models.}
  \label{fig:optimism}
\end{figure}

A gradient-based optimiser might exploit model error, finding plans whose predicted return the environment does not
deliver. We measure each executed plan's optimism: its predicted return, whose terminal value carries the goal bonus,
minus the discounted return actually received from that step on, using steps at least 299 steps before the episode
ends so that the part of the future the episode did not reach weighs under 5\%. Every planner is optimistic by about
2, as expected of a value fitted to the 0.7 expectile, but by nearly the same amount (Fig.~\ref{fig:optimism}): 2.06
(pure SCP) to 2.19 at $N=10$, 2.20 to 2.44 at $N=20$ and 2.41 to 2.56 at $N=30$, with SCP+MPPI inside the samplers'
range at every horizon and no trend with the sampling budget. We find no evidence that SCP, or a larger sampling
budget, exploits the model more than sampling with 512 sequences does.

% ---------------------------------------------------------------------------------------------------------------------
\subsection{Warned and unwarned violations}\label{app:crashes}

\begin{figure}[htbp]
  \centering
  \includegraphics[width=\linewidth]{\figdir/ablation/fig_crash_types.pdf}
  \caption{Violations per 100 episodes over the six $N=10$ models and $\epsilon\ge0.1$, split into those preceded
  within $N$ steps by a plan that broke the barrier in the model (warned, light) and those that were not (unwarned,
  dark, counts above the bars).}
  \label{fig:crash_types}
\end{figure}

Few violations come without warning (Fig.~\ref{fig:crash_types}). SCP+MPPI has 14 unwarned violations in 880
episodes, against 6 and 7 for the 512-sequence samplers and 7 and 13 for the time-matched ones. That stronger
optimisation, by SCP or by more samples, leads to more violations the model did not foresee would be consistent with
it planning closer to where the model is wrong, but paired by seed none of these differences is significant
($p\ge0.12$).

% ---------------------------------------------------------------------------------------------------------------------
\subsection{When violations happen}\label{app:survival}

\begin{figure}[htbp]
  \centering
  \includegraphics[width=\linewidth]{\figdir/appendix/fig_survival_pg1_c1_5.pdf}
  \caption{Share of episodes without a violation by each step, at each margin, on PointGoal1's $c=1.5$ model at
  $N=10$. Episodes end early only at a violation, so these are Kaplan--Meier curves without censoring.}
  \label{fig:survival}
\end{figure}

Most violations happen within the first 300 or so steps, after which the curves flatten (Fig.~\ref{fig:survival});
pure SCP's keeps falling for longer.
 in the appendix. Files and numbers as in experiments_results.tex.
% Order is by importance; drop subsections from the end if the appendix must be short.
% =====================================================================================================================
\providecommand{\figdir}{../paper_graphs}

\section{Additional Results}\label{app:results}

% ---------------------------------------------------------------------------------------------------------------------
\subsection{Per-model results}\label{app:per_model}

\begin{table}[htbp]
  \centering
  \caption{Table~\ref{tab:pooled} on the $c=1.2$ world models.}
  \label{tab:pooled_c1_2}
  \small
  \input{\figdir/appendix/table_pooled_c1_2}
\end{table}

\begin{table}[htbp]
  \centering
  \caption{Table~\ref{tab:pooled} on the unbounded ($c=\infty$) world models.}
  \label{tab:pooled_cinf}
  \small
  \input{\figdir/appendix/table_pooled_cinf}
\end{table}

Table~\ref{tab:per_model} gives each planner's violation rate and return on every $N=10$ world model, pooled over
$\epsilon\ge0.1$, and Figs.~\ref{fig:grid_violations} and~\ref{fig:grid_return} show the full sweeps. Pure SCP is
significantly less safe than SCP+MPPI on every model. Against the samplers, the only significant difference is MPPI
with 512 sequences on PointGoal2 with $c=1.2$, which violates less.

\begin{table}[htbp]
  \centering
  \caption{Each planner on each $N=10$ world model, pooled over $\epsilon\ge0.1$: violation rate (Wilson) and return
  (bootstrap) with 95\% intervals, and McNemar's $p$ for SCP+MPPI against it on paired seeds, Holm-corrected over the
  six models.}
  \label{tab:per_model}
  \small
  \input{\figdir/appendix/table_per_model}
\end{table}

\begin{figure}[htbp]
  \centering
  \includegraphics[width=\linewidth]{\figdir/appendix/fig_eps_sweep_all_bounds_violations.pdf}
  \caption{Violation rate against $\epsilon$ for every planner on every $N=10$ world model (rows: tasks; columns:
  Lipschitz bound $c$), with 95\% Wilson intervals.}
  \label{fig:grid_violations}
\end{figure}

\begin{figure}[htbp]
  \centering
  \includegraphics[width=\linewidth]{\figdir/appendix/fig_eps_sweep_all_bounds_return.pdf}
  \caption{Return against $\epsilon$ for every planner on every $N=10$ world model, with 95\% bootstrap intervals.}
  \label{fig:grid_return}
\end{figure}

% ---------------------------------------------------------------------------------------------------------------------
\subsection{Lipschitz bound and SCP's convergence}\label{app:lipschitz}

\begin{figure}[htbp]
  \centering
  \includegraphics[width=\linewidth]{\figdir/appendix/fig_lipschitz_convergence.pdf}
  \caption{Share of SCP+MPPI's steps that reached the 40-iteration cap against $\epsilon$, for each Lipschitz bound,
  at $N=10$.}
  \label{fig:lipschitz_convergence}
\end{figure}

Without a Lipschitz bound, SCP+MPPI reaches the iteration cap on more steps than with either bound at every margin,
except at $\epsilon=0.5$ on PointGoal2, where $c=1.5$ and $c=\infty$ are within 0.02 (Fig.~\ref{fig:lipschitz_convergence}).
Which bound converges best differs by task: $c=1.5$ on PointGoal1, $c=1.2$ on PointGoal2.

% ---------------------------------------------------------------------------------------------------------------------
\subsection{What more samples buy at \texorpdfstring{$N=10$}{N=10}}\label{app:samples}

\begin{figure}[htbp]
  \centering
  \includegraphics[width=\linewidth]{\figdir/ablation/fig_sample_vs_time_matched.pdf}
  \caption{Fig.~\ref{fig:eps_sweep} with the samplers at both budgets: 512 sequences (open markers, dashed) and
  time-matched (filled, solid).}
  \label{fig:sample_vs_time_matched}
\end{figure}

\begin{figure}[htbp]
  \centering
  \includegraphics[width=\linewidth]{\figdir/ablation/fig_return_vs_violations.pdf}
  \caption{Return against violation rate on each task's $c=1.5$ model at $N=10$, pooled over $\epsilon\ge0.1$, with
  95\% intervals. Arrows lead from each sampler with 512 sequences to the same sampler time-matched.}
  \label{fig:trade_off}
\end{figure}

Figs.~\ref{fig:sample_vs_time_matched} and~\ref{fig:trade_off} show what 128 times as many sequences buy the samplers.
On PointGoal1's $c=1.5$ model, the time-matched samplers earn more than with 512 sequences but violate more often
(MPPI: return 0.78 to 1.15, violation rate 0.19 to 0.24; CEM: 0.84 to 0.94 and 0.14 to 0.26), while SCP+MPPI reaches
the higher return at the lower violation rate (1.15 and 0.17). On PointGoal2 the extra sequences raise return (0.35 to
0.51 and 0.49) without changing the violation rate. These shifts lie within the intervals; the paired tests of
Sec.~\ref{sec:q1} find no significant difference in safety between SCP+MPPI and either budget.

% ---------------------------------------------------------------------------------------------------------------------
\subsection{All budgets at every horizon}\label{app:horizon_budgets}

\begin{figure}[htbp]
  \centering
  \includegraphics[width=\linewidth]{\figdir/ablation/fig_horizon_all_budgets.pdf}
  \caption{Fig.~\ref{fig:horizon} with the samplers at both 512 sequences (open markers, dashed) and their
  time-matched budget (filled, solid).}
  \label{fig:horizon_all}
\end{figure}

With 512 sequences the samplers' return at $N=20$ is close to SCP+MPPI's (MPPI 2.07, CEM 2.58, against 2.34; neither
difference significant), and at $N=30$ below it (1.59 and 1.22 against 2.26; $p=0.043$ and $0.013$)
(Fig.~\ref{fig:horizon_all}). Together with Fig.~\ref{fig:compute}, this shows that the samplers' best budget depends on
the sampler and the horizon.

% ---------------------------------------------------------------------------------------------------------------------
\subsection{Does optimisation exploit the model?}\label{app:optimism}

\begin{figure}[htbp]
  \centering
  \includegraphics[width=\linewidth]{\figdir/ablation/fig_model_optimism.pdf}
  \caption{Model optimism: each executed plan's predicted return (discounted predicted rewards plus terminal value)
  minus the discounted return received from that step on, averaged over each episode, with 95\% bootstrap intervals.
  $N=10$ pools every world model ($\epsilon\ge0.1$); $N=20$ and $30$ are the horizon study's models.}
  \label{fig:optimism}
\end{figure}

A gradient-based optimiser might exploit model error, finding plans whose predicted return the environment does not
deliver. We measure each executed plan's optimism: its predicted return, whose terminal value carries the goal bonus,
minus the discounted return actually received from that step on, using steps at least 299 steps before the episode
ends so that the part of the future the episode did not reach weighs under 5\%. Every planner is optimistic by about
2, as expected of a value fitted to the 0.7 expectile, but by nearly the same amount (Fig.~\ref{fig:optimism}): 2.06
(pure SCP) to 2.19 at $N=10$, 2.20 to 2.44 at $N=20$ and 2.41 to 2.56 at $N=30$, with SCP+MPPI inside the samplers'
range at every horizon and no trend with the sampling budget. We find no evidence that SCP, or a larger sampling
budget, exploits the model more than sampling with 512 sequences does.

% ---------------------------------------------------------------------------------------------------------------------
\subsection{Warned and unwarned violations}\label{app:crashes}

\begin{figure}[htbp]
  \centering
  \includegraphics[width=\linewidth]{\figdir/ablation/fig_crash_types.pdf}
  \caption{Violations per 100 episodes over the six $N=10$ models and $\epsilon\ge0.1$, split into those preceded
  within $N$ steps by a plan that broke the barrier in the model (warned, light) and those that were not (unwarned,
  dark, counts above the bars).}
  \label{fig:crash_types}
\end{figure}

Few violations come without warning (Fig.~\ref{fig:crash_types}). SCP+MPPI has 14 unwarned violations in 880
episodes, against 6 and 7 for the 512-sequence samplers and 7 and 13 for the time-matched ones. That stronger
optimisation, by SCP or by more samples, leads to more violations the model did not foresee would be consistent with
it planning closer to where the model is wrong, but paired by seed none of these differences is significant
($p\ge0.12$).

% ---------------------------------------------------------------------------------------------------------------------
\subsection{When violations happen}\label{app:survival}

\begin{figure}[htbp]
  \centering
  \includegraphics[width=\linewidth]{\figdir/appendix/fig_survival_pg1_c1_5.pdf}
  \caption{Share of episodes without a violation by each step, at each margin, on PointGoal1's $c=1.5$ model at
  $N=10$. Episodes end early only at a violation, so these are Kaplan--Meier curves without censoring.}
  \label{fig:survival}
\end{figure}

Most violations happen within the first 300 or so steps, after which the curves flatten (Fig.~\ref{fig:survival});
pure SCP's keeps falling for longer.
. Citations: experiments_refs.bib.
%
% Every figure and table is written by the repository's scripts, run from the repository root:
%   python src/scripts/analyse_results.py  -> src/save_data/experiments/analysis_results/*.csv
%   python src/scripts/paper_figures.py    -> paper_graphs/{main_text,ablation,appendix}/ (figures and both tables)
%   python src/scripts/env_figure.py       -> paper_graphs/main_text/fig_environments.pdf
% \figdir is the paper_graphs folder; define it before \input if the paper lives elsewhere (e.g. on Overleaf).
%
% Notation follows the method: N horizon, \epsilon barrier margin, \alpha DCBF rate, c Lipschitz bound, \tau risk
% threshold, g_\Phi margin head. Labels of the method's parts used here (rename to match yours): sec:regularity,
% sec:cbf, sec:scp, sec:init, eq:ocp.
% Every number is from the analysis of the runs on main (2d4240b). "% TODO" marks a fact to confirm before submission.
% =====================================================================================================================
\providecommand{\figdir}{../paper_graphs}

\section{Experimental Setup}\label{sec:setup}

Our experiments answer four questions.
\textbf{(Q1)}~Does SCP, started from a sampled plan, plan better than sampling-based MPC given the same wall-clock time?
\textbf{(Q2)}~How does the comparison change as the horizon, and with it the number of decision variables, grows?
\textbf{(Q3)}~What does SCP's convergence tell us online?
\textbf{(Q4)}~Does meeting the barrier in the learned model keep the robot safe in the environment, and what do the
regularity constraints cost?

% ---------------------------------------------------------------------------------------------------------------------
\subsection{Environments}\label{sec:envs}

\begin{figure}[t]
  \centering
  \includegraphics[width=0.62\linewidth]{\figdir/main_text/fig_environments.pdf}
  \caption{The two tasks from above, one layout each (layouts are drawn from the episode's seed). The point robot
  (red) must reach the goal (green) without entering a hazard (blue); in PointGoal2, touching a vase (cyan) is also a
  violation.}
  \label{fig:envs}
\end{figure}

We use \texttt{SafetyPointGoal1} and \texttt{SafetyPointGoal2} from Safety-Gymnasium~\citep{ji2023safety}
(Fig.~\ref{fig:envs}). A point robot with two actuators (forward force and turning, $u\in[-1,1]^2$) must reach a goal
of radius $0.3$\,m, which moves to a new position each time it is reached, within $T=1000$ steps. The reward is the
per-step reduction in distance to the goal plus a bonus of 1 for reaching it. The robot observes 60 quantities: 12
proprioceptive signals (accelerometer, velocimeter, gyroscope and magnetometer) and three 16-ray pseudo-lidars measuring
the proximity of the goal, the hazards and the vases. Level~1 places 8 hazards of radius $0.2$\,m and one cost-free
vase in a $3\times3$\,m region; level~2 places 10 hazards and 10 vases in a $4\times4$\,m region and also charges for
touching or moving a vase. A step is a \emph{violation} if it incurs any cost. We treat violations as absorbing
failures: an episode ends at its first violation, so the \emph{violation rate} is the share of episodes that fail and
the \emph{return} is the reward collected before failure or the time limit.

These tasks suit our question for four reasons.
(i)~DSRL~\citep{liu2023datasets} provides offline datasets for both, so the world model is learned without any
interaction, which is the setting of this paper.
(ii)~The constraint is geometric (stay out of regions of the plane) and observed through the lidar, so the barrier's
margin is a distance in metres, the margin $\epsilon$ has a physical meaning, and safety in the model can be checked
against safety in the environment (Sec.~\ref{sec:latent_safety}).
(iii)~Level~2 is more cluttered (20 obstacles in 16\,m$^2$ against 8 in 9\,m$^2$) and adds a second, movable kind of
obstacle, so the two levels, together with $\epsilon$, vary how restrictive the safe set is.
(iv)~The action is two-dimensional, which favours sampling: at $N=10$ a plan has only 20 decision variables, so the
comparison is a demanding one for SCP. The horizon study (Sec.~\ref{sec:horizon}) raises this to 40 and 60.

% ---------------------------------------------------------------------------------------------------------------------
\subsection{Offline data and world models}\label{sec:wm_setup}

For each task we train on the corresponding DSRL dataset (2,022 trajectories for PointGoal1 and 3,442 for PointGoal2,
of up to 1000 steps each), truncating every trajectory at its first violation to match the absorbing failure above.
This leaves about 0.87\,M transitions for PointGoal1 and 1.15\,M for PointGoal2; only 329 of PointGoal2's trajectories
run their full length without a violation.
% TODO: PG1's exact transition count (FilterDataset.ipynb prints it for PG2 only: 1,145,452); one sentence on how DSRL
% collected the data (behaviour policies, cost budgets).
The encoder drops seven proprioceptive channels (the accelerometer and the out-of-plane components of velocity,
angular rate and magnetic field) and maps the remaining 53 to a latent state of $d_z=32$ dimensions. The four networks
the planner differentiates through (dynamics, reward, value and margin) are MLPs with two hidden layers of 256 units,
each with Lipschitz constant at most $c$ (Sec.~\ref{sec:regularity}). We train one model per task and bound,
$c\in\{1.2,1.5,\infty\}$ ($c=\infty$: unbounded), for $N=10$, and two more PointGoal1 models with $c=1.5$ for $N=20$
and $30$. Each is trained for 250 epochs of 1000 (PointGoal1) or 1500 (PointGoal2) minibatches of 256 windows of $N+1$
consecutive transitions, with learning rate $3\times10^{-4}$, the loss at rollout step $t$ weighted by $\rho^t$
($\rho=0.9$, $0.95$, $0.99$ for $N=10$, $20$, $30$), discount $\gamma=0.99$ and value expectile $0.7$. The margin head
$g_\Phi$ regresses the distance from the robot to the nearest hazard edge (in PointGoal2, hazard or vase edge),
computed from the lidar, clipped at 1\,m, with samples near an edge up-weighted.
% TODO: optimiser (Adam?); state that each model is a single training seed (limitations).

% ---------------------------------------------------------------------------------------------------------------------
\subsection{Planners}\label{sec:planners}

All planners solve the same receding-horizon problem~\eqref{eq:ocp} in the same model: maximise
$\sum_{k=0}^{N-1}\gamma^k r(z_k,u_k)+\gamma^N V(z_N)-\rho_u\|u\|^2$ ($\rho_u=10^{-4}$) subject to
$(1-\alpha)h(z_k)-h(z_{k+1})\le\tau=0$ at every step, with $h(z)=g_\Phi(z)-\epsilon$ and $\alpha=0.05$. Each executes
the first action of its plan and replans at every step.

\textbf{SCP} is the method of Sec.~\ref{sec:scp} started from its previous plan shifted by one step (zeros at the first
step): at most 40 prox-linear iterations, each a convex QP solved by Clarabel~\citep{goulart2024clarabel}, with exact
$\ell_1$ penalties $w_{\mathrm{dyn}}=10^3$ and $w_{\mathrm{safe}}=10^2$, a trust-region weight that starts at 1, is
floored at $10^{-2}$ and is carried between steps, acceptance thresholds $(0.25,0.75)$, and a stop when the
prox-gradient norm falls below $10^{-2}$.
\textbf{SCP+MPPI}, our method, first runs MPPI with 512 sequences (below) and starts SCP from whichever of MPPI's plan
and its own shifted plan has the lower penalty merit (Sec.~\ref{sec:init}). SCP accepts only steps that lower the
merit, so the plan SCP+MPPI executes is never worse under the merit than MPPI's, whether or not SCP converges.
\textbf{MPPI}~\citep{williams2017information} and \textbf{CEM}~\citep{rubinstein1999cross} run 6 iterations of $J$
action sequences drawn from a Gaussian that is refitted to the best $J/8$ (weighted by $\exp(\text{score}/0.5)$ in MPPI,
uniformly in CEM) and warm-started from the previous mean shifted by one step. A sequence whose rollout breaks the
barrier at any step ranks below every sequence that meets it; if none meets it, sequences are ranked by their total
barrier excess, the quantity SCP's $\ell_1$ penalty charges. The executed plan is the refitted mean if its rollout
meets the barrier, and otherwise the best sequence sampled.

We run the samplers with $J=512$, the proposal SCP+MPPI uses, and \emph{time-matched}, with $J$ chosen so that a step
takes about as long as SCP+MPPI's: $J=65{,}536$ at $N=10$ (about 0.70\,s against 0.40--0.71\,s for SCP+MPPI, timed on
the same idle GPU), and $95{,}040$ and $120{,}000$ at $N=20$ and $30$ (mean logged plan times 2.04\,s and 3.81\,s,
against 2.16\,s and 3.61\,s for SCP+MPPI). Each run used one 1g.10gb slice of an NVIDIA A100, with the QP solved on the
CPU.
% TODO: the 65,536 time is interpolated between 32,768 (0.33 s) and 131,072 (1.43 s) sequences; time it directly or
% say so. Confirm the hardware sentence.

% ---------------------------------------------------------------------------------------------------------------------
\subsection{Protocol and statistics}\label{sec:protocol}

We sweep the margin $\epsilon\in\{0,0.1,0.2,0.3,0.5\}$\,m on every $N=10$ model and fix $\epsilon=0.2$\,m for the
horizon study. Each configuration runs on seeds $0$--$39$, or $0$--$29$ for the time-matched samplers, the horizon
study, PointGoal2's $c=\infty$ model and SCP+MPPI on PointGoal1's $c=\infty$ model. The seed fixes the layout, the
initial state and the planner's random numbers, so every comparison between planners or models is paired by seed. We
report violation rates with Wilson 95\% intervals and mean returns with bootstrap 95\% intervals ($10^4$ resamples).
Paired comparisons report the difference in violation rate and in return with a 95\% bootstrap interval, and the
probability, given the data, that SCP+MPPI's violation rate is lower or its return higher than the baseline's, from a
Bayesian bootstrap~\citep{rubin1981bayesian}. When pooling over margins and tasks we use $\epsilon\ge0.1$, since at
$\epsilon=0$ every planner fails often (Sec.~\ref{sec:q1}), and resample whole seeds, since a seed's episodes share a
layout. The per-model results in the appendix use McNemar's exact test on violations, Holm-corrected over the six models,
and the Wilcoxon signed-rank test on returns.

% =====================================================================================================================
\section{Results}\label{sec:results}

% ---------------------------------------------------------------------------------------------------------------------
\subsection{SCP needs a sampled start, and then matches time-matched sampling at \texorpdfstring{$N=10$}{N=10}}\label{sec:q1}

\begin{figure}[t]
  \centering
  \includegraphics[width=\linewidth]{\figdir/main_text/fig_eps_sweep_c1_5.pdf}
  \caption{Violation rate (top) and return (bottom) against the margin $\epsilon$ on each task's $c=1.5$ model at
  $N=10$, with 95\% intervals; 40 episodes per point (30 for the time-matched samplers).}
  \label{fig:eps_sweep}
\end{figure}

\begin{table}[t]
  \centering
  \caption{SCP+MPPI minus each baseline on the $c=1.5$ world models at $N=10$, pooled over both tasks and
  $\epsilon\ge0.1$ and paired by task, margin and seed (Tables~\ref{tab:pooled_c1_2} and~\ref{tab:pooled_cinf} in the
  appendix give the other bounds). A negative $\Delta$ violation rate or a positive $\Delta$ return favours SCP+MPPI.
  Brackets: 95\% intervals resampling whole seeds. $P$(safer) and $P$(higher return): the probability, given the data,
  that SCP+MPPI's violation rate is lower, or its return higher, than the baseline's on these models (Bayesian bootstrap
  over seeds).}
  \label{tab:pooled}
  \small
  \input{\figdir/main_text/table_pooled_c1_5}
\end{table}

\textbf{SCP needs a sampled start.} Started from its own shifted plan, SCP fails in 55\% of episodes over all $N=10$
models and $\epsilon\ge0.1$, against 25\% for SCP+MPPI on the same layouts (Table~\ref{tab:pooled}: $-29.8$ points,
$p=2\times10^{-43}$), at the same return. Pure SCP does not fail for want of convergence: on PointGoal1's $c=1.5$ model
it certifies nearly as many steps as SCP+MPPI (62--95\% against 63--97\% across $\epsilon$; Sec.~\ref{sec:certificates}). It fails
as a local method does, converging to the local solution nearest a poor start: it converges with barrier slack in use,
meaning no safe plan nearby, on more steps (2--19\% against 0--11\% on PointGoal1) and finds a plan that meets the
barrier in the model less often (88\% of steps against 92\%). On the $c=1.5$ models, SCP+MPPI starts from MPPI's plan on
33--43\% of steps on PointGoal1 and 38--55\% on PointGoal2 ($\epsilon\ge0.1$).

\textbf{With a sampled start, SCP+MPPI is as safe as sampling and earns more than its proposal.} Against MPPI and CEM
with the 512 sequences that seed it, SCP+MPPI's violation rate differs by $+1.9$ and $+2.6$ points ($p=0.24$ and
$0.098$), while its return is higher by 0.21 and 0.18 with intervals that exclude zero. Against samplers given the same
time, with 128 times as many sequences, the differences in violations ($-0.1$ and $+0.7$ points) and return ($+0.06$
and $+0.09$) are small and not significant. At $N=10$, then, SCP's refinement adds about as much return as 128 times
more samples, and neither makes the planner safer: with a two-dimensional action, sampling finds plans that meet the
barrier about as often as optimisation does (Sec.~\ref{sec:latent_safety}). The per-model results agree
(Table~\ref{tab:per_model}): no time-matched sampler differs significantly from SCP+MPPI in violations on any model,
and the one significant difference against the 512-sequence samplers favours MPPI on PointGoal2 with $c=1.2$ (23\%
against 34\%, Holm-corrected $p=0.031$).

\textbf{The margin is the main safety lever.} Without a margin ($\epsilon=0$) every planner fails in 30--63\% of
episodes, and pure SCP in 72--85\% (Fig.~\ref{fig:eps_sweep}): the barrier holds only in the model, so a margin must
absorb the model's error. A margin of 0.1--0.3\,m cuts SCP+MPPI's violation rate from 40--50\% to 15--23\% at a moderate
cost in return. At $\epsilon=0.5$\,m the planners turn conservative and return collapses (SCP+MPPI on PointGoal1: 0.38
against 1.75 at $\epsilon=0.1$), and on the more cluttered PointGoal2 violations rise again for every planner
(40--45\%; pure SCP 70\%) as the barrier becomes infeasible in many states, which SCP reports (Sec.~\ref{sec:certificates}).

% ---------------------------------------------------------------------------------------------------------------------
\subsection{At longer horizons SCP+MPPI earns more than time-matched sampling}\label{sec:horizon}

\begin{figure}[t]
  \centering
  \includegraphics[width=\linewidth]{\figdir/main_text/fig_horizon.pdf}
  \caption{Violation rate and return against the horizon $N$ on PointGoal1 ($c=1.5$, $\epsilon=0.2$), each sampler at
  its horizon's time-matched budget, with 95\% intervals. Each $N$ uses a world model trained on windows of $N+1$
  steps; 30 episodes per point (40 for SCP+MPPI and SCP at $N=10$).}
  \label{fig:horizon}
\end{figure}

% CUT 1 if over the page limit: move this figure to the appendix and keep the paragraph's numbers.
\begin{figure}[t]
  \centering
  \includegraphics[width=\linewidth]{\figdir/main_text/fig_return_vs_compute.pdf}
  \caption{Return against mean plan time per step at $N=20$ and $30$, with 95\% intervals. Each sampler is shown at
  every budget we ran (sequences per iteration along the bottom); SCP+MPPI and SCP are single points.}
  \label{fig:compute}
\end{figure}

We repeat the comparison on PointGoal1 ($c=1.5$, $\epsilon=0.2$) at $N=20$ and $30$, re-matching the samplers' budgets
in time at each horizon (Fig.~\ref{fig:horizon}). At both horizons SCP+MPPI earns more than the time-matched samplers,
at the same safety (2 and 1 violations in 30 episodes, against 3 and 0). Its return is higher by 0.65 [0.03, 1.30]
(MPPI, $p=0.021$) and 0.93 [0.15, 1.69] (CEM, $p=0.009$) at $N=20$, and by 1.13 [0.41, 1.93] ($p<0.001$) and 0.98
[0.08, 1.93] ($p=0.005$) at $N=30$.

The samplers do not turn extra time into return (Fig.~\ref{fig:compute}). Their largest budgets earn less than smaller
ones at both horizons, and their best budget changes with the sampler and the horizon (512 sequences for CEM at
$N=20$, 32,768 at $N=30$). At its best budget each sampler ties with SCP+MPPI (CEM: $-0.24$ at $N=20$, $+0.11$ at
$N=30$; neither significant), but that budget is not known in advance. SCP+MPPI used the same settings at every
horizon, and its refinement adds 0.27 (not significant) and 0.67 ($p=0.043$) to the return of the 512-sequence MPPI
plan it starts from.
% Possible reason for the samplers' trend, not tested: the elite set grows with J (top-k = J/8), so a larger budget
% refits to the same top eighth of the samples instead of concentrating on the best of them.

Two caveats apply. The horizon study uses one world model per horizon, one margin and 30 seeds, and the model changes
with $N$, so returns compare within a horizon but not across horizons. And within its 40-iteration cap SCP+MPPI
converged on 53\% of steps at $N=20$ and 32\% at $N=30$, executing its last accepted iterate on the rest: the
comparison is at matched time, not at convergence.

% ---------------------------------------------------------------------------------------------------------------------
\subsection{SCP's convergence certifies the plan and detects infeasibility}\label{sec:certificates}

\begin{figure}[t]
  \centering
  \includegraphics[width=\linewidth]{\figdir/main_text/fig_scp_stop_reasons.pdf}
  \caption{How SCP+MPPI's steps ended on each task's $c=1.5$ model at $N=10$: converged with the plan's rollout meeting
  the barrier (certified), converged with barrier slack in use, or at the 40-iteration cap.}
  \label{fig:stops}
\end{figure}

Each SCP step ends in one of three ways (Fig.~\ref{fig:stops}). Because SCP's reference trajectory is always the model
rollout of its plan, a step that stops as converged, with prox-gradient norm below $10^{-2}$ and barrier excess below
$10^{-4}$, executes an approximately stationary plan that meets the barrier at every step in the model: convergence
itself certifies the plan, at no extra cost. This is the most common outcome. On the $c=1.5$ models SCP+MPPI certifies
97\% of steps at $\epsilon=0$, falling to 63\% at $\epsilon=0.5$, on PointGoal1, and 83\% falling to 60\% on
PointGoal2, using 9--22 iterations on average.

A step that converges with barrier slack in use signals something different: a stationary point at which the barrier
cannot be met, that is, no safe plan near the current one. Its share grows with the margin, from 0 to 11\% of steps on
PointGoal1 and from 1\% to 22\% on PointGoal2, where at $\epsilon=0.5$ it coincides with the rise in violations
(Sec.~\ref{sec:q1}). Sampling planners give no such signal: when no sampled sequence meets the barrier they cannot tell
an empty safe set from too few samples. The remaining steps reach the iteration cap, a share that grows with $\epsilon$
(3\% to 27\% on PointGoal1) and with the horizon (46\% at $N=20$ and 68\% at $N=30$ on the same task and model bound).

% ---------------------------------------------------------------------------------------------------------------------
\subsection{Safety in the model predicts safety in the environment}\label{sec:latent_safety}

% CUT 2 if over the page limit: move this figure to the appendix; the paragraph stands on its numbers.
\begin{figure}[t]
  \centering
  \includegraphics[width=\linewidth]{\figdir/main_text/fig_latent_safety.pdf}
  \caption{Share of steps followed by a violation within $N$ steps, after a plan whose model rollout met the barrier
  and after one whose rollout did not, over the six $N=10$ models and $\epsilon\ge0.1$, with 95\% Wilson intervals
  (log scale).}
  \label{fig:latent_safety}
\end{figure}

The barrier is enforced only in the learned model; Fig.~\ref{fig:latent_safety} tests whether that is enough. After a
plan that meets the barrier in the model, a violation follows within $N$ steps on $2$--$10\times10^{-4}$ of steps,
whichever planner made the plan (including the 512-sequence samplers); after a plan that does not, on
$3$--$7\times10^{-2}$, 56 to 131 times as often. Conversely, 93--98\% of violations were preceded within $N$ steps by a
plan that broke the barrier in the model: the model nearly always warns before a violation. Violations therefore arise
mainly in states where the planner finds no plan that is safe in the model, not from plans that are safe in the model
but not in the environment. All planners find such a plan on a similar share of steps (90--92\%; pure SCP 88\%), which
is why optimising harder does not make the planner safer at $N=10$ (Sec.~\ref{sec:q1}).

% ---------------------------------------------------------------------------------------------------------------------
\subsection{Regularity constraints}\label{sec:regularity_results}

Bounding the networks' Lipschitz constants at $c=1.5$ costs no measurable safety or return: against the unbounded
models, no planner's violation rate or return differs significantly on either task (McNemar $p\ge0.18$; return
intervals include zero;
Figs.~\ref{fig:grid_violations} and~\ref{fig:grid_return} in the appendix). It does make SCP converge more reliably
(Fig.~\ref{fig:lipschitz_convergence}). At $\epsilon$ from 0.1 to 0.3, SCP+MPPI reaches the iteration cap on 6--16\% of
steps on PointGoal1 and 16--21\% on PointGoal2 with $c=1.5$, against 33--44\% and 31--35\% unbounded, needs fewer
iterations (12--16 against 24--28 on PointGoal1), and so certifies more of its plans (78--92\% against 55--65\% on
PointGoal1; 74--83\% against 53--62\% on PointGoal2). The tighter bound $c=1.2$ makes the model less conservative:
returns are higher, but so are violations, significantly on PointGoal1 for SCP+MPPI (44 against 28 of 160 episodes,
$p=0.044$), MPPI ($p=0.009$) and CEM ($p=0.005$), and on PointGoal2 for SCP+MPPI (54 against 38, $p=0.036$). We
therefore use $c=1.5$ for the main results.
